\documentclass[10pt]{amsart}
\usepackage[a4paper,margin=23mm]{geometry}
\usepackage{amsmath,amssymb,amsthm,xfrac,hyperref,xspace}
\usepackage[english]{babel}
\usepackage[scr=rsfs,cal=euler]{mathalfa}
\newcommand{\resp}{resp.}
\newcommand{\cf}{cf.}
\newcommand{\ie}{i.e.}
\newcommand{\eg}{e.g.}
\newcommand{\loccit}{loc.\ cit.}
\newcommand{\skpt}{\par\smallskip\noindent}
\let\Subsection\subsection
\let\Subsubsection\subsubsection
\emergencystretch=3em
\multlinegap0pt
\hyphenation{belong-ing defined denote denotes series}
\let\ts\textstyle

\newcommand{\Changel}{\mathcode`l="7160}
\newcommand{\Changelback}{\mathcode`l="716C}
\newcommand{\NoChangel}[1]{%
\expandafter\let\csname old\string#1\endcsname=#1
\let#1=\relax
\newcommand{#1}{\mathop{\mathcode`l="716C\csname old\string#1\endcsname\mathcode`l="7160}\displaylimits}%
}
\newcommand{\NoChangelnolimits}[1]{%
\expandafter\let\csname old\string#1\endcsname=#1
\let#1=\relax
\newcommand{#1}{\mathop{\mathcode`l="716C\csname old\string#1\endcsname\mathcode`l="7160}\nolimits}%
}

\NoChangelnolimits{\log}
\NoChangelnolimits{\ln}
\NoChangel{\lim}
\NoChangel{\limsup}
\NoChangel{\liminf}
\NoChangel{\varlimsup}
\NoChangel{\varliminf}
\NoChangel{\varinjlim}
\NoChangel{\varprojlim}

\newenvironment{smallbmatrix}{\big[\begin{smallmatrix}}{\end{smallmatrix}\big]}

\newenvironment{Smallbmatrix}{\Big[\begin{smallmatrix}}{\end{smallmatrix}\Big]}

\newcommand{\Psfrac}[2]{(\sfrac{#1}{#2})}
\newcommand{\psfrac}[2]{\sfrac{(#1)}{#2}}
\newcommand{\spfrac}[2]{\sfrac{#1}{(#2)}}
\newcommand{\pspfrac}[2]{\sfrac{(#1)}{(#2)}}
\newcommand{\Ppsfrac}[2]{(\sfrac{(#1)}{#2})}

\newcommand\mto{\mathchoice{\longmapsto}{\mapsto}{\mapsto}{\mapsto}}
\newcommand\hto{\mathrel{\lhook\joinrel\to}}
\newcommand\To[1]{\mathchoice{\xrightarrow{\textstyle\kern4pt#1\kern3pt}}{\stackrel{#1}{\longrightarrow}}{}{}}

\newcommand{\Lfunction}{$L$\nobreakdash-func\-tion\xspace}
\newcommand{\Lfunctions}{$L$\nobreakdash-func\-tions\xspace}

\usepackage{bm}
\usepackage{dsfont}
\usepackage{bbold}
\usepackage{leftidx}
\usepackage{stmaryrd}
\usepackage{tikz-cd}

\newenvironment{bsm}
{\left[\begin{smallmatrix}}
{\end{smallmatrix}\right]}

\theoremstyle{plain}
\newtheorem{thm}{Theorem}[subsection]
\newtheorem{prop}[thm]{Proposition}
\newtheorem{thmintro}{Theorem}[section]

\theoremstyle{remark}
\newtheorem{rmk}[thm]{Remark}
\newtheorem{rmkintro}[thmintro]{Remark}

\let\mr\mathrm

\newcommand{\bA}{\mathbb{A}}
\newcommand{\bC}{\mathbb{C}}
\newcommand{\bI}{\mathbb{I}}
\newcommand{\bQ}{\mathbb{Q}}
\newcommand{\bS}{\mathbb{S}}
\newcommand{\bT}{\mathbb{T}}
\newcommand{\bU}{\mathbb{U}}
\newcommand{\bZ}{\mathbb{Z}}

\newcommand{\bba}{\mathbb{a}}
\newcommand{\bbb}{\mathbb{b}}
\newcommand{\bbc}{\mathbb{c}}

\newcommand{\cB}{\mathcal{B}}
\newcommand{\cC}{\mathcal{C}}
\newcommand{\cD}{\mathcal{D}}
\newcommand{\cE}{\mathcal{E}}
\newcommand{\cF}{\mathcal{F}}
\newcommand{\cH}{\mathcal{H}}
\newcommand{\cK}{\mathcal{K}}
\newcommand{\cL}{\mathcal{L}}
\newcommand{\cM}{\mathcal{M}}
\newcommand{\cN}{\mathcal{N}}
\newcommand{\cO}{\mathcal{O}}
\newcommand{\cP}{\mathcal{P}}
\newcommand{\cS}{\mathcal{S}}
\newcommand{\cW}{\mathcal{W}}

\newcommand{\fa}{\mathfrak{a}}
\newcommand{\fb}{\mathfrak{b}}
\newcommand{\fc}{\mathfrak{c}}
\newcommand{\fd}{\mathfrak{d}}
\newcommand{\fp}{\mathfrak{p}}
\newcommand{\fx}{\mathfrak{x}}
\newcommand{\fz}{\mathfrak{z}}

\newcommand{\fA}{\mathfrak{A}}
\newcommand{\fB}{\mathfrak{B}}
\newcommand{\fD}{\mathfrak{D}}

\newcommand{\sB}{\mathscr{B}}
\newcommand{\sC}{\mathscr{C}}
\newcommand{\sS}{\mathscr{S}}
\newcommand{\sT}{\mathscr{T}}

\DeclareMathAlphabet{\mathpzc}{OT1}{qzc}{m}{it}

\newcommand{\pzV}{\mathpzc{V}}

\newcommand{\bfP}{\mathbf{P}}

\newcommand{\bbeta}{\bm{\beta}}

\newcommand{\utau}{\protect\underline{\tau}}

\DeclareMathOperator{\GL}{GL}
\DeclareMathOperator{\GSp}{GSp}
\DeclareMathOperator{\GU}{GU}
\DeclareMathOperator{\U}{U}
\DeclareMathOperator{\SL}{SL}
\DeclareMathOperator{\Sp}{Sp}

\newcommand{\Meas}{\mathpzc{M}\mkern-2mu eas}

\newcommand{\cont}{\mathrm{cont}}
\newcommand{\diag}{\mathrm{diag}}
\newcommand{\Nm}{\mathrm{Nm}}
\newcommand{\ord}{\mathrm{ord}}
\newcommand{\Tr}{\mathrm{Tr}}
\newcommand{\triv}{\mathrm{triv}}
\newcommand{\tor}{\mathrm{tor}}
\newcommand{\val}{\mathrm{val}}\NoChangel{\val}
\newcommand{\vol}{\mathrm{vol}}\NoChangel{\vol}
\newcommand{\Sieg}{\mr{Sieg}}
\newcommand{\bdot}{\boldsymbol{\cdot}}
\newcommand{\disc}{\mathrm{disc}}
\newcommand{\spe}{\mathrm{sp}}
\newcommand{\KL}{\mathrm{KL}}

\DeclareMathOperator{\Gal}{Gal}\NoChangel{\Gal}
\DeclareMathOperator{\Her}{Her}
\DeclareMathOperator{\Hom}{Hom}
\DeclareMathOperator{\Spec}{Spec}
\DeclareMathOperator{\Sym}{Sym}

\newcommand{\adic}{\text{-}\mathrm{adic}}
\newcommand{\qexp}{q\text{-}\mathrm{exp}}

\let\ol\overline
\let\wh\widehat
\let\wt\widetilde
\let\mf f
\let\Bes\cB
\let\bes B
\let\Schw\Phi
\let\Whi\cW
\let\whi W
\let\llb\llbracket
\let\rrb\rrbracket

\newcommand{\bid}{\mathbf{1}}

\newcommand{\dsec}{{\texttt f}}

\newcommand{\alphaS}{\alpha_{\mathbb{S}}}
\newcommand{\alphabS}{\bar{\alpha}_{\mathbb{S}}}
\newcommand{\sPi}{\scriptscriptstyle{\Pi}}
\newcommand{\sLambda}{\scriptscriptstyle{\Lambda}}
\newcommand{\bLam}{\boldsymbol\Lambda}
\newcommand{\bUps}{\boldsymbol\Upsilon}


\begin{document}
\title{$p$-adic $L$-functions for $\mathrm{GSp}(4)\times\mathrm{GL}(2)$, II}
\author{Zheng Liu}
\maketitle
\setcounter{section}{1}
\begin{thmintro}\label{thm:main}
Given Hida families $\sC_1,\sC_2$ on $\GL(2),\GSp(4)$ and the auxiliary data:
\begin{itemize}
\item
$\bS=\begin{bmatrix}\bba&\sfrac{\bbb}{2}\\ \sfrac{\bbb}{2}&\bbc\end{bmatrix}\in\Sym_2(\bQ)_{>0}$,
\item
a $p$-adically continuous Hecke character $\bLam:\cK^\times\backslash\bA^\times_{\cK,f}\to \Lambda^\times_{\GSp(4)}$ with $\cK=\bQ(\sqrt{-\det\bS})$, such that $\bLam|_{\bA^\times_\bQ}=\omega_{\sC_2}$, the central character associated to $\sC_2$,
\item
a finite set $S$ of places of $\bQ$ containing $p,\infty$ such that everything is unramified outside $S$, (see Section~\ref{sec:setup} for the precise condition on $S$),
\item
an open subgroup $U^p$ of $\hat{\bZ}^{p,\times}=\prod_{v\neq \infty,p} \bZ^\times_v$ containing $\prod_{v\notin S} \bZ^\times_v$,
\end{itemize}
taking $\beta_1\in\bQ_{>0},\beta_2\in\Sym_2(\bQ)_{>0}$, there exists a four-variable $p$-adic \Lfunction
\begin{multline*}
\cL^S_{\sC_1,\sC_2,\beta_1,\beta_2}\in\Meas\bigl(\bQ^\times\backslash\bA^\times_{\bQ,f}/U^p,\bI_{\sC_1}\wh{\otimes}\bI_{\sC_2}\bigr)\otimes_{\bI_{\sC_1}\wh{\otimes}\bI_{\sC_2}} (F_{\sC_1}\wh{\otimes}F_{\sC_2})\\
\cong (\bI_{\sC_1}\wh{\otimes}\bI_{\sC_2})\llb \bQ^\times\backslash\bA^\times_{\bQ,f}/U^p\rrb\otimes_{\bI_{\sC_1}\wh{\otimes}\bI_{\sC_2}} (F_{\sC_1}\wh{\otimes}F_{\sC_2})
\end{multline*}
\noindent
\begin{minipage}{\textwidth}
satisfying the interpolation property:
\begin{multline*}
\cL^S_{\sC_1,\sC_2,\beta_1,\beta_2}(\kappa,x)
=2^{-l-l_1-l_2}i^l\sum_{f\in \sS_{\GL(2),x}} \frac{f_{\bbc}f_{\beta_1}}{\bfP(f,f)}\sum_{\varphi\in \sS_{\GSp(4)},x} \frac{\bes^\dagger_{\bS,\Lambda}\left(\varphi\right)\varphi_{\beta_2}}{\bfP(\varphi,\varphi)} \\
\times E_\infty \Bigl(k+\frac{l+l_1+l_2}{2},\Pi_x\times\pi_x\times \chi\Bigr) E_p\Bigl(k+\frac{l+l_1+l_2}{2},\Pi_x\times\pi_x\times \chi\Bigr)\\
\times L^S\Bigl(k+\frac{l+l_1+l_2}{2},\Pi_x\times\pi_x\times \chi\Bigr),
\end{multline*}
\end{minipage}
\noindent
where
\begin{itemize}
\item
$x\in \sC_1(\bar{\bQ}_p)\times\sC_2(\bar{\bQ}_p)$ is a point at which the weight projection map \eqref{eq:wtproj} is \'{e}tale with arithmetic image
\[
(\tau,(\tau_1,\tau_2))=((l,\xi),(l_1,l_2,\xi_1,\xi_2))
\in \Hom_\cont\Bigl(T^1_{\GL(2)}(\bZ_p)\times T^1_{\GSp(4)}(\bZ_p),\bar{\bQ}^\times_p\Bigr),
\]
and $\kappa=(k,\chi)$ is an arithmetic point in $\Hom_\cont\bigl(\bQ^\times\backslash\bA^\times_{\bQ,f}/U^p,\bar{\bQ}^\times_p\bigr)$, (see Section~\ref{sec:pchar} for some of the notations), such that
\begin{equation}
\begin{aligned}\label{eq:kl}
-\frac{\min\{-l_1+l_2+l,\,l_1+l_2-l\}}{2}+2 &\leq k+\frac{l_1+l_2+l}{2}\\
&\leq \frac{\min\{-l_1+l_2+l,\,l_1+l_2-l\}}{2}-1,
\end{aligned}
\end{equation}
(when $\min\{-l_1+l_2+l,\,l_1+l_2-l\}\geq 3$, $s=k+\psfrac{l+l_1+l_2}{2}$ for such $k$'s are all the critical points for the \Lfunction $L(s,\Pi_x\times\pi_x\times\chi)$),
\item
$\sS_{\GL(2),x}$ (\resp $\sS_{\GSp(4),x})$ is an orthogonal basis (with respect to the modified \hbox{Petersson} inner product defined in \eqref{eq:Pet2}, \eqref{eq:Pet}) of the space spanned by ordinary cuspidal holomorphic forms on $\GL(2)$ of weight $l$ and tame level $K^p_{\GL(2)}$ (\resp $\GSp(4)$ of weight $(l_1,l_2)$ and tame level $K^p_{\GSp(4)}$) with nebentypus at $p$ given by \eqref{eq:pneben}, belonging to the Hecke eigenspace parameterized by $x$,
\item
$\pi_x$ (\resp $\Pi_x$) is any unitary cuspidal irreducible automorphic representation of $\GL(2,\bA_\bQ)$ (\resp $\GSp(4,\bA_\bQ)$) inside the representation generated by $\sS_{\GL(2),x}$ (\resp $\sS_{\GSp(4),x}$) twisted by a real power of $|\det|$,
\item
the factors $E_p,E_\infty$ are the modified Euler factors for $p$-adic interpolation (see \eqref{eq:Ep}, \eqref{eq:Einf} for the precise formula),
\item
$f_\bbc,f_{\beta_1}$ (\resp $\varphi_{\beta_2}$) denote the Fourier coefficient of $f$ indexed by $\bbc,\beta_1$ (\resp of $\varphi$ indexed by $\beta_2$), the modified Petersson inner products $\bfP(f,f),\bfP(\varphi,\varphi)$ are defined as in \eqref{eq:Pet2}, \eqref{eq:Pet},
\item
$\Lambda$ is the classical Hecke character corresponding to the specialization of $\bLam$ at $(\tau_1,\tau_2)$, and $\bes^\dagger_{\bS,\Lambda}(\varphi)$ is the Bessel period with a modification at $p$. (See \eqref{eq:bes-dagger} and \cite[\S2.2.1]{pFL} for the precise definition.)
\end{itemize}
\end{thmintro}
\section{Notation and review of Hida theory and Furusawa's formula}
\numberwithin{equation}{subsection}

\subsection{Notation}

We fix an odd prime number $p$, an isomorphism $\bar{\bQ}_p\cong\bC$, and a sufficiently large finite extension $F$ of $\bQ_p$. Denote by $\cO$ the ring of integers of $F$.

We use $v$ to denote a place of $\bQ$. We~fix the additive character
\[
\psi_{\bA_\bQ}=\bigotimes_v\psi_v:\bQ\backslash\bA\to\bC^\times,
\qquad \psi_v(x)=
\begin{cases}
e^{-2\pi i\{x\}_v}&\text{if } v\neq\infty,
\\e^{2\pi ix}&\text{if } v=\infty,
\end{cases}
\]
where $\{x\}_v$ is the fractional part of $x$, \ie $\{x\}_v\in\bQ\cap[0,1)$ with $x-\{x\}_v\in\bZ_v$.

Given a positive integer $n$, denote by $\bid_n$ the identity matrix of size $n$, and define the algebraic group $\GSp(2n)$ over $\bZ$ as
\begin{align}
\setlength{\arraycolsep}{2.5pt}
\label{eq:GSp2n}\GSp(2n,R)= \Bigl\{g\in\GL(2n,R):\ltrans{g}\begin{smallbmatrix}0&\bid_n\\-\bid_n&0\end{smallbmatrix}g=\nu_g\begin{smallbmatrix}0&\bid_n\\-\bid_n&0\end{smallbmatrix}, \,\nu_g\in R^\times \Bigr\}
\end{align}
for all $\bZ$-algebra $R$. Given an imaginary quadratic field $\cK$, define the algebraic group $\GU(n,n)$ over $\bZ$ as
\begin{multline}
\setlength{\arraycolsep}{2.5pt}
\label{eq:GUnn}\GU(n,n)(R)\\
=\Bigl\{g\in\GL(2n,\cO_\cK\otimes R):\ltrans{\bar{g}}\begin{smallbmatrix}0&\bid_n\\-\bid_n&0\end{smallbmatrix}g=\nu_g\begin{smallbmatrix}0&\bid_n\\-\bid_n&0\end{smallbmatrix}, \,\nu_g\in R^\times \Bigr\},
\end{multline}
where for $\alpha\in\cK$, $\bar{\alpha}$ denotes its image under the nontrivial element in $\Gal(\cK/\bQ)$. Given $g\in \GSp(2n)$ or $\GU(n,n)$, we~call the $\nu_g$ in \eqref{eq:GSp2n}\eqref{eq:GUnn} the similitude of $g$. In this paper, we~will work with $\GSp(4),\GSp(2)=\GL(2),\GU(3,3),\GU(1,1)$.

Fix the following maximal torus of $\GSp(2n)$, $\GU(n,n)$:
\begin{align*}
T_{\GSp(2n)}&=\bigl\{\diag(a_1,\dots,a_n,\nu a^{-1}_1,\dots,\nu a^{-1}_n)\in\GSp(2n)\bigr\},\\
T_{\GU(n,n)}&=\bigl\{\diag(\fa_1,\dots,\fa_n,\nu \bar{\fa}^{-1}_1,\dots,\nu \bar{\fa}^{-1}_n)\in\GU(n,n)\bigr\},
\end{align*}
and Siegel parabolic subgroups
\begin{equation}\label{eq:Q}
\begin{aligned}
Q_{\GSp(2n)}&=\Bigl\{\begin{Smallbmatrix} A&B\\0&\nu \ltrans{A}^{-1}\end{Smallbmatrix}\in \GSp(2n) \Bigr\},\\
Q_{\GU(n,n)}&= \Bigl\{\begin{Smallbmatrix} \fA&\fB\\0&\nu \ltrans{\bar{\fA}}^{-1}\end{Smallbmatrix}\in \GU(n,n) \Bigr\},
\end{aligned}
\end{equation}
where the matrices are written in $n\times n$ blocks. Denote by $M_{\GSp(2n)}\subset Q_{\GSp(2n)}$, $M_{\GU(n,n)}\subset Q_{\GU(n,n)}$ the Levi subgroups. Let
\begin{align*}
T^1_{\GSp(2n)}&=T_{\GSp(2n)}\cap \Sp(2n),
&T^1_{\GU(n,n)}&=T_{\GU(n,n)}\cap \U(n,n),\\
M^1_{\GSp(2n)}&=M_{\GSp(2n)}\cap \Sp(2n),
&M^1_{\GU(n,n)}&=M_{\GU(n,n)}\cap \U(n,n),
\end{align*}
where $\Sp(2n)$ (\resp $\U(n,n)$) is the subgroup of $\GSp(2n)$ (\resp $\GU(n,n)$) consisting of elements with similitude $1$. We~identify $M^1_{\GSp(2n)},M^1_{\GU(n,n)}$ with $\GL(n),\GL(n)_{/\cK}$~via
\[
A\mto \begin{bmatrix} A\\&\ltrans{A}^{-1}\end{bmatrix},\qquad
\fA\mto \begin{bmatrix} \fA\\&\ltrans{\bar{\fA}}^{-1}\end{bmatrix}.
\]
Denote by
\[
U_{M^1_{\GSp(2n)}}\subset M^1_{\GU(n,n)},\qquad
U_{M^1_{\GU(n,n)}}\subset M^1_{\GU(n,n)}
\]
the subgroup whose elements are upper triangular with diagonal entries being $1$ under the above identification. Let
\begin{align*}
U_{\GSp(2n)}&= \Bigl\{\begin{Smallbmatrix} A&B\\0&\nu \ltrans{A}^{-1}\end{Smallbmatrix}\in Q_{\GSp(2n)}:A\in U_{M^1_{\GSp(2n)}} \Bigr\},\\
U_{\GU(n,n)}&= \Bigl\{\begin{Smallbmatrix} \fA&\fB\\0&\nu \ltrans{\bar{\fA}}^{-1}\end{Smallbmatrix}\in Q_{\GU(n,n)}:\fA\in U_{M^1_{\GU(n,n}} \Bigr\}.
\end{align*}
For all $\bZ$-algebras $R$, we~use $\Her_n(\cO_\cK\otimes R)$ (\resp $\Sym_n(R)$) to denote the set of all $n\times n$ matrices $\sigma$ with entries in $\cO_\cK\otimes R$ (\resp $R$) satisfying $\ltrans{\bar{\sigma}}=\sigma$ (\resp $\ltrans{\sigma}=\sigma$).

Given a character $\theta:\bQ^\times_p\to \bC^\times$ or a character $\Theta:\cK^\times_p\to\bC^\times$, we~let
\begin{equation}
\label{eq:char-circ}
\theta^\circ=\mathds{1}_{\bZ^\times_p}\cdot \theta,
\qquad
\Theta^\circ=\mathds{1}_{\cO^\times_{\cK,p}}\cdot \Theta.
\end{equation}
Given a Hecke character $\Theta:\cK^\times\backslash\bA^\times_\cK\to\bC^\times$, we~let
\[
\Theta_\bQ=\Theta|_{\bA^\times_\bQ}.
\]
We will use $\triv$ to denote the trivial character.

\subsection{Review of Hida theory}\label{sec:Hida}
We recall some constructions and results in Hida theory for symplectic groups. (We will use it for $\GSp(4)$ and $\GSp(2)=\GL(2)$.) See \cite{HidaCon}or \cite[\S6.1]{SLF} for details.

For $G=\GSp(2n)$, define the Iwasawa algebra
\begin{align}
\label{eq:Iw-alg}
\wt{\Lambda}_G=\cO\llb T^1_G(\bZ_p)\rrb,
\qquad
\Lambda_G=\cO\llb T^1_G(1+p\bZ_p)\rrb\cong\cO\llb T_1,T_2,\dots,T_n\rrb.
\end{align}
Fix a neat open compact subgroup $K^p_G\subset G(\bA^p_{\bQ,f})$. Let $Y_G$ denote the Shimura variety for $G$ of level $K^p_G G(\bZ_p)$ defined over $\cO$. Let $\sT_{G,l,m}$ denote the $l$-th layer of the Igusa tower over $\bZ/p^m\bZ$, which is an $M^1_G(\bZ/p^l)$-\'{e}tale cover of the ordinary locus $Y_{G,\ord}$, and $\sT^\tor_{G,l,m}$ be a smooth partial toroidal compactification of $\sT_{G,l,m}$ (with respect to a chosen polyhedral cone decomposition) with boundary $C$. (See \cite[\S\S2.1,\,3.1,\,3.2]{HidaCon}.) Put
\[
V_{G,l,m}=H^0\left(\sT^\tor_{G,l,m},\cO_{\sT_{G,l,m}}(-C)\right)^{U_{M^1_G(\bZ_p)}},
\]
and
\begin{align*}
V_G=\varprojlim_m\varinjlim_l V_{G,l,m},
\qquad\pzV_G=\varinjlim_m\varinjlim_l V_{G,l,m}.
\end{align*}
The group $T^1_G(\bZ_p)$ naturally acts on these spaces, and they are all naturally modules over $\wt{\Lambda}_G$ and $\Lambda_G$.

For a tuple of integers $\underline{m}:m_1\!\geq\! m_2\!\geq\!\cdots\!\geq\! m_n\!\geq\! 0$ and $m_0\!\geq\! 0$, an~oper\-ator $U^G_{p,\underline{m},m_0}$ (corresponding to $\diag(p^{m_1+m_0},\dots,p^{m_n+m_0},p^{-m_1},\dots,p^{-m_n})$) on~$V_{G,l,m}$ is defined. We~call these operators $\bU_p$-operators. Put $U^G_p=U^G_{p,\underline{m},0}$ with $\underline{m}=(n,n-\nobreak1,\dots,1)$. As operators on $V_{G,l,m}$, the limit
\[
e^G_\ord=\lim_{r\to \infty} \left(U^G_p\right)^{r!}
\]
exists \cite[Th.\,1.1 (3)]{HidaCon}\cite[\S6.1.3]{SLF}, and is called the ordinary projector. Let
\begin{align*}
V_{G,\ord}&= e^G_\ord V_G,
&\pzV_{G,\ord}&= e^G_\ord \pzV_G,
&\pzV^*_{G,\ord}&=\Hom(\pzV_{G,\ord},F/\cO).
\end{align*}
The $\tilde{\Lambda}_G$-module of Hida families of cuspidal $p$-adic automorphic forms on $G$ of tame level $K^p_G$ is defined to be
\[
\cM_{G,\ord}=\Hom_{\wt{\Lambda}_G}(\pzV^*_G,\wt{\Lambda}_G).
\]
For each $p$-adic weight $\underline{\tau}\in \Hom_{\cont}\bigl(T^1_G(\bZ_p),\bar{\bQ}^\times_p\bigr)$, there is a natural map
\begin{multline}\label{eq:embd}
V_{G,\ord}[\utau]
\to \Hom_{\cO}\bigl(\Hom(V_{G,\ord},\cO)/\cP_{\utau},\wt{\Lambda}_G/\cP_{\utau}\bigr)\\
\to \Hom_{\cO}\bigl(\Hom(\pzV_{G,\ord},F/\cO)/\cP_{\utau},\wt{\Lambda}_G/\cP_{\utau}\bigr) \to \cM_{G,\ord}\otimes_{\wt{\Lambda}_G} \wt{\Lambda}_G/\cP_{\utau},
\end{multline}
where $\cP_{\utau}$ is the prime ideal of $\wt{\Lambda}_G$ corresponding to $\utau$.

\begin{thm}[{\cite[Th.\,1.1]{HidaCon}}]
$\cM_{G,\ord}$ is free over $\Lambda_{G}$ of finite rank. The map \eqref{eq:embd} is an isomorphism. If~$\utau$ is algebraic and sufficiently regular, then
\[
V_{G,\ord}[\utau]=e^G_\ord H^0\bigl(Y^\tor_G,\omega_{\utau}(-C)\bigr),
\]
where $\omega_{\utau}$ is the automorphic vector bundle of weight $\utau$ over a toroidal compactification of $Y_G$. (The right hand side can be identified with the space of classical holomorphic automorphic forms on $G$ of weight $\utau$.)
\end{thm}

\subsection{Review of Furusawa's formula}\label{sec:Furusawa}
We quickly recall a modification of Furusawa's formula for \Lfunctions for $\GSp(4)\times\GL(2)$. See \cite[\S2.1]{pFL} for details. Take
\begin{align*}
\bS=\begin{bmatrix}\bba&\sfrac{\bbb}{2}\\\sfrac{\bbb}{2}&\bbc\end{bmatrix}\in\Sym_2(\bQ)_{>0},
\end{align*}
and let $\cK=\bQ(\sqrt{-\det \bS})$, and $\eta_{\cK/\bQ}:\bQ^\times\backslash\bA^\times_\bQ\to\bC^\times$ be the quadratic character corresponding to $\cK/\bQ$. Let
\begin{align*}
\alphaS=\frac{\bbb+\sqrt{\bbb^2-4\bba\bbc}}{2\bbc},\arraycolsep3.5pt
\qquad\imath_\bS\left(\fz\right)
=\begin{bmatrix}\alphaS&1\\ \alphabS&1\end{bmatrix}^{-1} \begin{bmatrix}\fz\\&\bar{\fz}\end{bmatrix}\begin{bmatrix}\alphaS&1\\ \alphabS&1\end{bmatrix}
\end{align*}
for $\fz\in \cK\otimes_\bQ R$ with $R$ any $\bQ$-algebra.

Given a Hecke character $\Xi:\cK^\times\backslash\bA^\times_\cK\to\bC^\times$, denote by $I_v(s,\chi,\Xi)$ the degenerate principal series on $\GU(3,3)(\bQ_v)$ consisting of smooth functions $\dsec_v(s,\chi,\Xi):\GU(3,3)(\bQ_v)\to\bC$ such that
\begin{align*}
\dsec_v(s,\chi,\Xi)\left(\begin{Smallbmatrix}\fA&\fB\\0&\fD\end{Smallbmatrix}g\right)=\Xi_v(\det \fA)\chi_v(\det \fA\fD^{-1})|\det \fA\fD^{-1}|^{s+\sfrac{3}{2}}_v\,\dsec_v(s,\chi,\Xi)(g)
\end{align*}
for all $g\in\GU(3,3)(\bQ_v)$ and $\begin{smallbmatrix}\fA&\fB\\0&\fD\end{smallbmatrix}\in Q_{\GU(3,3)}(\bQ_v)$, the Siegel parabolic group defined in \eqref{eq:Q}. The Siegel Eisenstein series associated to a section $\dsec(s,\chi,\Xi)\in I(s,\chi,\Xi)=\bigotimes'_v I_v(s,\chi,\Xi)$ is defined as
\[
E^{\Sieg}(g;\dsec(s,\chi,\Xi))=\sum_{\gamma\in Q_{\GU(3,3)}(\bQ)\backslash \GU(3,3)(\bQ)}\dsec(s,\chi,\Xi)(\gamma g).
\]

Let $\pi$ be an irreducible cuspidal automorphic representation of $\GL(2,\bA_\bQ)$. By taking a Hecke character $\Upsilon:\cK^\times\backslash\bA^\times_\cK\to\bC^\times$ with $\Upsilon_\bQ$ equal to the central character of $\pi$, for every $f\in\pi$, we~can extend it to an automorphic form $f^\Upsilon$ on $\GU(1,1)$ by
\begin{equation}\label{eq:extGL2}
\mf^\Upsilon(\fa g)=\Upsilon(\fa)f(g), \qquad\fa\in\bA^\times_\cK,\,g\in\GL(2,\bA_\bQ).
\end{equation}
Then $\pi^{\Upsilon}=\{f^\Upsilon:f\in\pi\}$ is an irreducible cuspidal automorphic representation of $\GU(1,1)$. Denote the Whittaker period of $f\in \pi$ with respect to $\psi_{\bA_\bQ,\bbc}$ (defined as $\psi_{\bA_\bQ,\bbc}(x)=\psi_{\bA_\bQ}(\bbc x)$) by $\whi_\bbc(f)$, and define the function $\Whi_\bbc(f)$ on $\GL(2,\bA_\bQ)$ as
\[
\Whi_\bbc(f)(g)=\whi_\bbc(g\cdot f).
\]

Let $\Pi$ be an irreducible cuspidal automorphic representation of $\GSp(4,\bA_\bQ)$, and $\Lambda:\cK^\times\backslash\bA^\times_\cK\to\bC^\times$ be a Hecke character such that $\Lambda_\bQ$ equals the central character of $\Pi$. Then for $\varphi\in \Pi$, one can define its global Bessel period $\bes_{\bS,\Lambda}(\varphi)$ with respect to $\bS,\Lambda$ (and $\psi$) as recalled in \cite[\S2.2.1]{pFL}. We~also define the function $\Bes_{\bS,\Lambda}(\varphi)$ on $\GSp(4,\bA_\bQ)$ as
\begin{align*}
\Bes_{\bS,\Lambda}(\varphi)(g)=\bes_{\bS,\Lambda}(g\cdot \varphi).
\end{align*}

Define $\GSp(4)\times_{\GL(1)}\GU(1,1)$ as $\left\{(g,h)\in \GSp(4)\times\GU(1,1):\nu_g=\nu_h\right\}$, and similarly define $\GU(2,2)\times_{\GL(1)}\GU(1,1)$. Denote by $\imath:\GU(2,2)\times_{\GL(1)}\GU(1,1)\hto\GU(3,3)$ the embedding
\[\arraycolsep4.5pt
\left(\begin{bmatrix}A&B\\C&D\end{bmatrix},\begin{bmatrix}\fa&\fb\\\fc&\fd\end{bmatrix}\right)\mto
\begin{bmatrix}A&&B\\&\fa&&\fb\\C&&D\\&\fc&&\fd\end{bmatrix}.
\]
Let $R'_\bS \subset \GSp(4)\times_{\GL(1)}\GU(1,1)$ be the subgroup
\[
\biggl\{\left(\begin{Smallbmatrix}\imath_\bS(\fz)\\&\ltrans{\imath_\bS(\bar{\fz})}\end{Smallbmatrix}\begin{Smallbmatrix}\bid_2&X\\&\bid_2\end{Smallbmatrix},\,\fz\cdot\bid_2\right):\fz\in \mr{Res}_{\cK/\bQ}\GL(1),\,X\in\Sym_2\biggr\}.
\]
Put
\[\arraycolsep4.5pt
\eta_\bS=\begin{bmatrix}
1\\ \alphaS&1\\&&1&-\alphabS\\&&&1
\end{bmatrix}\in\GU(2,2)(\bQ),
\qquad\cS=\begin{bmatrix}1\\&1\\&&1\\&&&1\\&&1&&1\\&1&&&&1\end{bmatrix}\in\GU(3,3)(\bQ).
\]

The following formula is proved in \cite[Th.\,2.1.1]{pFL} by combining Furusawa's formula \cite{Furusawa} and Garrett's generalization of the doubling method \cite{GaKl}.
\begin{thm}\label{thm:Furusawa}
Assume that $\Xi\Lambda^c\Upsilon^c=\triv$ and $S$ is a subset of places of $\bQ$ containing $\infty$ such that at all $v\notin S$,
\begin{itemize}
\item
$\pi,\Pi,\Upsilon,\Xi,\chi,\dsec(s,\chi,\Xi)$ are all unramified and $\varphi, f$ are invariant under the action by $\GSp(4,\bZ_v),\GU(1,1)(\bZ_v)$,
\item
$\bS=\begin{smallbmatrix}\bba&\sfrac{\bbb}{2}\\\sfrac{\bbb}{2}&\bbc\end{smallbmatrix}$ belongs to $M_2(\bZ_v)$ with $\bbc\in\bZ^\times_v$ and $\bbb^2-4\bba\bbc=\mr{disc}(\cK_v/\bQ_v)$.
\end{itemize}
\begin{minipage}{\textwidth}
Then
\begin{multline*}
\int_{[\GSp(4)\times_{\GL(1)}\GU(1,1)]} E^{\Sieg}\bigl(\imath(g,h);\dsec(s,\chi,\Xi)\bigr)\cdot \varphi(g) \cdot \mf^\Upsilon(h)\,\Xi^{-1}(\det h) \,dh\,dg\\
=\whi_\bbc(f)\cdot \bes_{\bS,\Lambda}(\varphi)\cdot d^S_3\Bigl(s+\frac{1}{2},\Xi(\chi\circ\Nm)\Bigr)^{-1}L^S\Bigl(s+\frac{1}{2},\tilde{\Pi}\times\tilde{\pi}\times\chi\Bigr)\\
\times \prod_{v\in S} Z_v\Bigl(\dsec_v(s,\chi,\Xi),\Bes^{\Pi_v}_{\bS,\Lambda_v}(\varphi_v), \Whi^{\pi_v,\Upsilon_v}_\bbc(f_v)\Bigr),
\end{multline*}
\end{minipage}
with
\[
d^S_3\bigl(s,\Xi(\chi\circ\Nm)\bigr)=\prod_{j=1}^3 L_S\bigl(2s+j,\Xi_\bQ\chi^2\eta^{n-j}_{\cK/\bQ}\bigr),
\]
and
\begin{multline*}
Z_v\Bigl(\dsec_v(s,\chi,\Xi),,\Bes^{\Pi_v}_{\bS,\Lambda_v}(\varphi_v), \Whi^{\pi_v,\Upsilon_v}_\bbc(f_v)\Bigr)
=\Bes^{\Pi_v}_{\bS,\Lambda_v}(\varphi_v)(\bid_4)^{-1}\Whi^{\pi_v,\Upsilon_v}_\bbc(f_v)(\bid_2)^{-1}\\
\times \int_{\bigl(R'_\bS\backslash\GSp(4)\times_{\GL(1)}\GU(1,1)\bigr)(\bQ_v)} \dsec_v(s,\chi,\Xi)\left(\cS^{-1}\imath(\eta_\bS\, g,h)\right)\\
\times \Bes^{\Pi_v}_{\bS,\Lambda_v}(\varphi_v)(g)
\Whi^{\pi_v,\Upsilon_v}_\bbc(f_v)\left(\begin{smallbmatrix}0&1\\-1&0\end{smallbmatrix}h\right)\Xi^{-1}_v(\det h)\,dhdg,
\end{multline*}
where $\Nm$ denotes the norm map from $\bA^\times_\cK$ to $\bA^\times_\bQ$, $\Bes^{\Pi_v}_{\bS,\Lambda_v}(\varphi_v)$ is the element corresponding to $\varphi$ in the local Bessel model of $\Pi_v$, $\Whi^{\pi_v,\Upsilon_v}_\bbc(f_v)$ is the extension to $\GU(1,1)(\bQ_v)$ via $\Upsilon_v$ of the element corresponding to $f$ in the local Whittaker model of $\pi_v$.
\end{thm}

In \cite{pFL}, a one-variable cyclotomic $p$-adic \Lfunction for $\Pi\times\pi$ is constructed by using the above integral and interpolating Siegel Eisenstein series when $s,\chi$ vary. (Such a modification of Furusawa's integral is also used in \cite{Morimoto-GSp4GL2-I,Morimoto-GSp4GL2-II}.) In the next section, we~let $\pi,\Pi$ also vary in Hida families, and construct a four-variable Siegel Eisenstein family.

\section{Four-variable \texorpdfstring{$p$-adic \protect\Lfunction for $\GSp(4)\times\GL(2)$}{pGL}}

\subsection{The setup}\label{sec:setup}
Let $S$ be a finite set of places of $\bQ$ containing $p,\infty$ and $U^p$ be an open subgroup of $\hat{\bZ}^{p,\times}$ containing $\prod_{v\notin S} \bZ^\times_v$. We~fix tame level groups\begin{align*}
\setlength{\arraycolsep}{2pt}
K^p_{\GL(2)}&=\prod_{v\notin S}\GL(2,\bZ_v)\times\prod_{v\in S-\{\infty\}} K_{\GL(2),v}, &&\text{with }\begin{Smallbmatrix}1&\bZ_v\\&1\end{Smallbmatrix}\subset K_{\GL(2),v},\\
K^p_{\GSp(4)}&=\prod_{v\notin S}\GSp(4,\bZ_v)\times\hspace{-.5em}\prod_{v\in S-\{\infty\}} K_{\GSp(4),v},&&\text{with }\begin{Smallbmatrix}\bid_2&\Sym_2(\bZ_v)\\&\bid_2\end{Smallbmatrix}\subset K_{\GSp(4),v}.
\end{align*}
Let $\bT_{\GL(2),\ord}$ (\resp $\bT_{\GSp(4),\ord}$) be the Hecke algebras acting on the $\tilde{\Lambda}_{\GL(2)}$-module $\cM_{\GL(2),\ord}$ of Hida families of tame level $K^p_{\GL(2)}$ (\resp $\tilde{\Lambda}_{\GSp(4)}$-module $\cM_{\GSp(4),\ord}$ of Hida families of tame level $K^p_{\GSp(4)}$) consisting of all the spherical Hecke operators away from $S$ and the $\bU_p$-operators at $p$. We~have the natural weight projection map
\begin{equation}\label{eq:wtproj}
\Spec(\bT_{\GL(2),\ord})\times\Spec(\bT_{\GSp(4),\ord})
\to \Spec(\tilde{\Lambda}_{\GL(2)})\times \Spec(\tilde{\Lambda}_{\GSp(4)}).
\end{equation}

\pagebreak[2]
Assume that we are given the following data:
\begin{itemize}
\item
a geometrically irreducible component $\sC_1\subset \Spec(\bT_{\GL(2),\ord})$ with central character
\[
\omega_{\sC_1}:\bQ^\times\backslash\bA^\times_{\bQ,f}\to \Lambda^\times_{\GL(2)},
\]
\item
a geometrically irreducible component $\sC_2\subset \Spec(\bT_{\GSp(4),\ord})$ with central character
\[
\omega_{\sC_2}:\bQ^\times\backslash\bA^\times_{\bQ,f}\to \Lambda^\times_{\GSp(4)},
\]
\end{itemize}
plus the auxiliary data:
\begin{itemize}
\item
an imaginary quadratic field $\cK$ and a positive definite symmetric form
\[\bS=\begin{bmatrix}\bba&\sfrac{\bbb}{2}\\ \sfrac{\bbb}{2}&\bbc\end{bmatrix}\in\Sym_2(\bQ)_{>0}
\]
with $\cK=\bQ(\sqrt{-\det\bS})$ such that for all $v\notin S$,
\begin{align*}
\bS\in \Sym_2(\bZ_v),
\qquad\bbc\in \bZ^\times_v,
\qquad 4\det\bS=\mr{disc}(\cK_v/\bQ_v)
\end{align*}
\item
a continuous character $\bUps:\cK^\times\backslash\bA^\times_{\cK,f}\to \Lambda^\times_{\GL(2)}$ extending $\omega_{\sC_1}$,
\item
a continuous character $\bLam:\cK^\times\backslash\bA^\times_{\cK,f}\to \Lambda^\times_{\GSp(4)}$ extending $\omega_{\sC_2}$.
\end{itemize}

\subsubsection{Notation for some \texorpdfstring{$p$}{p}-adic characters}\label{sec:pchar}
We identify $T^1_{\GL(2)}(\bZ_p), T^1_{\GSp(4)}(\bZ_p)$ with $\bZ^\times_p,\bZ^\times_p\times\bZ^\times_p$ via
\begin{align*}
a\mto \diag(a,a^{-1}),\qquad
(a_1,a_2)\mto \diag(a_1,a_2,a^{-1}_1,a^{-1}_2).
\end{align*}
We denote by $\tau$ (\resp $(\tau_1,\tau_2)$) a continuous character
\[
T^1_{\GL(2)}(\bZ_p)\to\bar{\bQ}^\times_p\qquad (\text{\resp} T^1_{\GSp(4)}(\bZ_p)\to\bar{\bQ}^\times_p). 
\]
When the characters are arithmetic, \ie
\begin{align*}
\tau(x)=x^l\xi(x),
\qquad
\tau_i(x)=x^{l_i}\xi_i(x),\quad i=1,2
\end{align*}
for some integers $l,l_1,l_2$ and finite order characters $\xi,\xi_1,\xi_2$ of $\bZ^\times_p$, we~write
\begin{align*}
\tau=(l,\xi),
\qquad
(\tau_1,\tau_2)=(l_1,l_2,\xi_1,\xi_2),
\end{align*}
and call $l$ (\resp $(l_1,l_2)$) the algebraic part of $\tau$ (\resp $(\tau_1,\tau_2)$), and $\xi$ (\resp $(\xi_1,\xi_2)$) the finite order part. We~view $\tau,\tau_1,\tau_2$ also as a $\bar{\bQ}_p$-valued character of $\bQ^\times\backslash \bA^\times_{\bQ,f}/\hat{\bZ}^{p,\times}$ via the isomorphism $\bZ^\times_p\To{\sim}\bQ^\times\backslash \bA^\times_{\bQ,f}/\hat{\bZ}^{p,\times}$ (induced by the embedding $\bZ^\times_p\hto \bA^\times_{\bA,f}$).

We denote by $\kappa$ a continuous character $\bQ^\times\backslash\bA^\times_{\bQ,f}/U^p \to\bar{\bQ}^\times_p$, and when it is arithmetic, we~write
\[
\kappa=(k,\chi)
\]
with $k$ the algebraic part and $\chi$ the finite order part.

In this paper, we~call an arithmetic tuple $(\kappa,\tau,\tau_1,\tau_2)$ classical if their algebraic parts satisfy
\begin{equation}\label{eq:k-crt}
-\min\{l_1+l_2,l+l_2\}+2\leq k\leq -\max\{l_1,l\}-1.
\end{equation}
This condition corresponds to
\[
\resizebox{\textwidth}{!}
{$\displaystyle
-\frac{\min\{-l_1+l_2+l,\,l_1+l_2-l\}}{2}+2 \leq k+\frac{l_1+l_2+l}{2}\leq \frac{\min\{-l_1+l_2+l,\,l_1+l_2-l\}}{2}-1
$}
\]
equivalent to $s=k+\psfrac{l_1+l_2+l}{2}$ being a critical point for the degree-8 \Lfunction $L(s,\Pi\times\pi)$ with $\Pi_\infty$ (\resp $\pi_\infty$) isomorphic to the holomorphic discrete series of weight $(l_1,l_2)$ (\resp $l$).

\subsubsection{Convention on nebentypus at \texorpdfstring{$p$}{p} and central characters}
Given arithmetic $\tau=(l,\xi)$, $(\tau_1,\tau_2)=(l_1,l_2,\xi_1,\xi_2)$, we~use the convention that the classical automorphic forms on $\GL(2),\GSp(4)$ contained in $V_{\GL(2)}[\tau],V_{\GSp(4)}[\tau_1,\tau_2]$ have
\begin{itemize}
\item
weight $l,(l_1,l_2)$,
\item
nebentypus such that the right translation of $B_{\GL(2)}(\bZ_p),B_{\GSp(4)}(\bZ_p)$ acts by the character
\begin{align}
\label{eq:pneben}
\arraycolsep4.5pt
\begin{bmatrix}a&\ast\\ &d\end{bmatrix}
\mto
\xi(a),
\qquad
\begin{bmatrix} a_1&\ast&\ast&\ast\\&a_2&\ast&\ast\\ && a^{-1}_1\nu\\ &&\ast& a^{-1}_2\nu\end{bmatrix}
\mto
\xi_1(a_1)\xi_2(a_2),
\end{align}
\item
central character equal to the product of a finite order character unramified away from $p$ and the classical character corresponding to
\begin{align*}
\tau:\bQ^\times\backslash \bA^\times_{\bQ,f}/\hat{\bZ}^{p,\times}\to \bar{\bQ}^\times_p,
\qquad\tau_1\tau_2:\bQ^\times\backslash \bA^\times_{\bQ,f}/\hat{\bZ}^{p,\times}\to \bar{\bQ}^\times_p.
\end{align*}
\end{itemize}
(Note that with this convention, the central characters of classical cuspidal automorphic forms in $V_{\GL(2)}[\tau],V_{\GSp(4)}[\tau_1,\tau_2]$ are not necessarily unitary.)

\Subsection{The Eisenstein measure}\label{sec:EisMeas}
\subsubsection{The Siegel Eisenstein series and its Fourier coefficients}
For a classical tuple $(\kappa,\tau,\tau_1,\tau_2)$ with $\kappa$ fixed by $U^p$, put
\begin{multline*}
E^\Sieg_{\kappa,\tau,\tau_1,\tau_2}(g)= |\nu_g|^{\frac{3}{2}(l_1+l_2+l)}
\cdot \frac{\prod_{j=0}^2\Gamma\left(s+\left|k+\frac{l+l_1+l_2}{2}-\frac{1}{2}\right|+3-j\right)}{2^{-6}(-2i)^{6\left(\left|k+\frac{l+l_1+l_2}{2}-\frac{1}{2}\right|+\frac{3}{2}\right)}\,\pi^{3\left(s+\left|k+\frac{l+l_1+l_2}{2}-\frac{1}{2}\right|+2\right)}}\\
\times d^S_3\bigl(s,\Lambda_0\Upsilon_0(\chi\circ\Nm)\bigr) \cdot
E^{\Sieg}\left(g;\dsec_{l,\,l_1,l_2,\xi,\,\xi_1,\xi_2}(s,\chi,\Lambda_0\Upsilon_0)\right)\Big|_{s=k+\frac{l_1+l_2+l}{2}-\frac{1}{2}},
\end{multline*}
where $\Lambda_0=\Lambda\,|\cdot|^{-\psfrac{l_1+l_2}{2}}$, $\Upsilon_0=\Upsilon \,|\cdot|^{-\sfrac{l}{2}}$ with $\Lambda,\Upsilon$ the classical characters corresponding to the specializations at $\tau_1,\tau_2,\tau$ of $\bLam,\bUps$ fixed in Section~\ref{sec:setup}, and the section
\[
\dsec_{l,\,l_1,l_2,\xi,\,\xi_1,\xi_2}(s,\chi,\Lambda_0\Upsilon_0)\in\bigotimes_v I_v(s,\chi,\Lambda_0\Upsilon_0)
\]
is chosen as in \cite[\S3.3]{pFL}. (For the place $p$, the $\eta^\circ_{\sPi_1,p},\eta^\circ_{\sPi_2,p},\eta^\circ_{\sPi_3,p}$ in \loccit correspond to $\xi^{-1}_2,\xi^{-1}_1,\xi^{-1}_1\xi^{-1}_2$ here, and $\eta^\circ_{\pi_p,1},\eta^\circ_{\pi_p,2}$ in \loccit correspond to $\xi^{-1},\triv$ here. For the Archimedean place, $k+\sfrac{\epsilon}{2}$ in \loccit correspond to $k+\psfrac{l_1+l_2+l}{2}$ here.)

Given $\bbeta\in\Her_3(\cK)$ and a nearly holomorphic automorphic form $F$ on $\GU(3,3)$, we~define the $\bbeta$-th polynomial Fourier coefficient of $F$ at $g_f\in \GU(3,3)(\bA_{\bQ,f})$ as the polynomial $F_{\bbeta}(g)\in \bC[Y]$, where $\bC[Y]$ denotes the ring of polynomials in entries of $Y=(Y_{ij})_{1\leq i,l\leq 3}$, such that for $z\in \left\{z\in M_{3,3}(\bC)\mid i(\ltrans{\bar{z}}-z)>0\right\}$,\vspace*{-5pt}
\begin{multline*}
F_{\bbeta}(g_f)\Bigl(Y=\Bigl(\frac{z-\ltrans{\bar{z}}}{2i}\Bigr)^{-1}\Bigr)\cdot e^{2\pi i\Tr z\bbeta}\\
=\int_{\Her_3(\cK)\backslash\Her_3(\bA_\cK)} F\biggl(\begin{bmatrix}\bid_3&\varsigma\\0&\bid_3\end{bmatrix}g_f\begin{bsm}\bigl(\tfrac{z-\ltrans{\bar{z}}}{2i}\bigr)^{1/2}&\tfrac{z+\ltrans{\bar{z}}}{2}\bigl(\tfrac{z-\ltrans{\bar{z}}}{2i}\bigr)^{-1/2}\\0&\bigl(\tfrac{z-\ltrans{\bar{z}}}{2i}\bigr)^{-1/2}\end{bsm}_\infty\biggr)\, \psi_{\bA_\bQ}\left(-\Tr\bbeta\,\fx\right)\,d\fx.
\end{multline*}
(The existence of such a polynomial $F_{\bbeta}(g)$ is implied by the definition of nearly holomorphic modular forms. See \cite[\S13.11]{Sh00} or \cite[\S3.2]{pFL} for the definition of nearly holomorphic forms and \cite[\S2.4]{NHF} for their interpretations as global sections of automorphic sheaves over Shimura varieties.) Similarly, one defines polynomial Fourier coefficients of nearly holomorphic forms on $\GU(1,1),\GL(2)$ (\resp $\GSp(4)$) indexed by $\bQ=\Her_1(\cK)$ (\resp $\Sym_2(\bQ)$).

By our choice of the Archimedean component of $ \dsec_{l,\,l_1,l_2,\xi,\,\xi_1,\xi_2}(s,\chi,\Lambda_0\Upsilon_0)$, the Eisenstein series $E^\Sieg_{\kappa,\tau,\tau_1,\tau_2}$ is a nearly holomorphic automorphic form on $\GU(3,3)$. We~have the following proposition on its polynomial Fourier coefficients.

\begin{prop}
For $\fA\in\GL(2,\bA^p_{\cK,f})$, $\nu\in \bA^{\times,p}_f$ and $\bbeta\in\Her_3(\cK)$, we have $\left(E^\Sieg_{\kappa,\tau,\tau_1,\tau_2}\right)_{\bbeta}=0 $ unless $\bbeta>0$, and for $\bbeta>0$,\vspace*{-3pt}
\[
\resizebox{\displaywidth}{!}{$\displaystyle \begin{aligned}
&\left(E^\Sieg_{\kappa,\tau,\tau_1,\tau_2}\right)_{\bbeta}\Bigl(\begin{Smallbmatrix}\fA\\&\nu\ltrans{\bar{\fA}}^{-1}\end{Smallbmatrix}^p_f\Bigr)(Y=0)
= A(\bbeta;\kappa,\tau,\tau_1,\tau_2) \Bigl(\begin{Smallbmatrix}\fA\\&\nu\ltrans{\bar{\fA}}^{-1}\end{Smallbmatrix}^p_f \Bigr)
\cdot C_{k,l,l_1,l_2}(\bbeta)
\end{aligned}$}\]
with\vspace*{-5pt}
\begin{multline}\label{eq:Abeta}
A(\bbeta;\kappa,\tau,\tau_1,\tau_2) \Bigl(\begin{Smallbmatrix}\fA\\\nu\ltrans{\bar{\fA}}^{-1}\end{Smallbmatrix}^p_f \Bigr)
=\Lambda\Upsilon(\det(\nu\ltrans{\bar{\fA}}^{-1}) \cdot (\chi|\cdot|^k)(\det(\nu\fA^{-1}\ltrans{\bar{\fA}}^{-1})) \\
\times\prod_{v\nmid Np\infty} h_{v,\nu^{-1}\ltrans{\bar{\fA}}\bbeta\fA}\left(\Lambda_{\bQ,v}\Upsilon_{\bQ,v}\chi^2_v(\varpi_v)|\varpi_v|^{2k+2}\right)
\prod_{v\mid N}\cF\Schw^{\vol}_{v,\val_v(N)} \bigl(\nu^{-1}_v\ltrans{\bar{\fA}}_v\bbeta\fA_v \bigr)\\
\times \Lambda^{\circ-1}_{\bar{\fp}}\xi_1(\varrho_\fp(\beta_{13}))\varrho_\fp(\beta_{13})^{-r_{\sLambda,2}+l_1} \cdot \Lambda^{\circ-1}_{\fp}\xi_1(\varrho_{\bar{\fp}}(\beta_{13}))\varrho_{\bar{\fp}}(\beta_{13})^{-r_{\sLambda,1}+l_1}\\
\times \xi_2\xi\chi^\circ_p \Bigl(\frac{\bar{\beta}_{13}\beta_{23}-\beta_{13}\bar{\beta}_{23}}{\alphaS-\alphabS}\Bigl)\Bigl(\frac{\bar{\beta}_{13}\beta_{23}-\beta_{13}\bar{\beta}_{23}}{\alphaS-\alphabS}\Bigr)^{l_2+l+k-2}\\
\times \xi_1\xi_2\chi^\circ_p\Bigl( \frac{(\alphaS-\alphabS)(\beta_{12}-\bar{\beta}_{12})}{2}\Bigr)\Bigl( \frac{(\alphaS-\alphabS)(\beta_{12}-\bar{\beta}_{12})}{2}\Bigr)^{l_1+l_2+k-2},
\end{multline}
and\vspace*{-5pt}
\begin{multline}\label{eq:Cbeta}
C_{k,l,l_1,l_2}(\bbeta)\\
=\begin{cases}
\biggl(\Bigl(\frac{\beta_{12}-\bar{\beta}_{12}}{2}\det\begin{Smallbmatrix}\bar{\beta}_{13}&\bar{\beta}_{23}\\ \beta_{13}&\beta_{23}\end{Smallbmatrix} \Bigr)^{-1}\det\bbeta) \biggr)^{2k+l_1+l_2+l-1} &\text{if }k+\psfrac{l_1+l_2+l}{2}\geq \sfrac{1}{2},\\
1&\text{if } k+\psfrac{l_1+l_2+l}{2}\leq \sfrac{1}{2}.
\end{cases}
\end{multline}
Here, $\Lambda,\Upsilon$ are the classical characters associated to the specialization of $\bLam$ at $(\tau_1,\tau_2)$, $\bUps$ at $\tau$, and $(r_{\sLambda,1},r_{\sLambda,2})$ is the $\infty$-type of $\Lambda$, $\varrho_p,\varrho_{\bar{\fp}}$ are embedding of $\cK$ into $\bar{\bQ}_p$ inducing $\fp,\bar{\fp}$. (See \eqref{eq:char-circ} for the notation superscript $^\circ$ on characters.) $\cF\Schw^{\vol}_{v,\val_v(N)}$ is the Fourier transform of the Schwartz function defined in \cite[\S3.3.3]{pFL}, which does not depend on $\kappa,\tau,\tau_1,\tau_2$.
\end{prop}

\begin{proof}
This proposition is the same as \cite[Prop.\,3.5.1]{pFL}, aside from a notational difference. The characters $\eta^\circ_{\sPi_1,p},\eta^\circ_{\sPi_2,p},\eta^\circ_{\sPi_3,p}$ in \loccit correspond to our $\xi^{-1}_2,\xi^{-1}_1,\xi^{-1}_1\xi^{-1}_2$ here, $\eta^\circ_{\pi_p,1},\eta^\circ_{\pi_p,2}$ in \loccit correspond to our $\xi^{-1},\triv$ here, and $k+\sfrac{\epsilon}{2},\Lambda,\Xi$ in \loccit correspond to our $k+\psfrac{l+l_1+l_2}{2}$, $\Lambda^{-c}_0,\Lambda_0\Upsilon_0$ here.
\end{proof}

\subsubsection{The \texorpdfstring{$p$}{p}-adic measure interpolating restrictions of Siegel Eisenstein series}
Given a compact and totally disconnected topological space $Y$ and a $p$-adically complete $\bZ_p$-module $M$, we~denote by $\Meas(Y,M)$ the space of all continuous $\bZ_p$-linear maps from $\sC(Y,\bZ_p)$ to $M$, where $\sC(Y,\bZ_p)$ is the space of all continuous $\bZ_p$-valued functions on $Y$ equipped with the topology of uniform convergence.
\begin{thm}\label{thm:EisF}
There exists a $p$-adic measure
\[
\mu_\cE\in\Meas\bigl(\bQ^\times\backslash\bA^\times_{\bQ,f}/U^p,\cM_{\GL(2),\ord}\wh{\otimes}_{\cO} \cM_{\GSp(4),\ord} \bigr)
\]
such that for all classical $(\kappa,\tau,\tau_1,\tau_2)=((k,\chi),(l,\xi),(l_1,l_2,\xi_1,\xi_2))$ with $k,l,l_1,l_2$ satisfying \eqref{eq:k-crt}, denoting by $\spe_{\tau,(\tau_1,\tau_2)}$ the specialization map at $\tau,(\tau_1,\tau_2)$, we~have
\begin{multline*}
\spe_{\tau,(\tau_1,\tau_2)}\left(\mu_\cE(\kappa)\right)
=e^{\GL(2)}_\ord e^{\GSp(4)}_\ord \,{\rm Proj}_{K^p_{\GL(2)},K^p_{\GSp(4)}} \bigl(E^\Sieg_{\kappa,\tau,\tau_1,\tau_2}\big|_{\GL(2)\times\GSp(4)} \bigr)\\
\times \spe_{\tau,(\tau_1,\tau_2)}(\omega^{-1}_{\sC_1}\omega^{-1}_{\sC_2}\circ{\rm det}_{\GL(2)}).
\end{multline*}
Here, $-\,|_{\GL(2)\times\GSp(4)}$ means restriction to $\GL(2)\times_{\GL(1)}\GSp(4)$ followed by extension by zero to $\GL(2)\times\GSp(4)$, and the projection ${\rm Proj}_{K^p_{\GL(2)},K^p_{\GSp(4)}}$ is defined as $\int_{K^p_{\GL(2)}}\int_{K^p_{\GSp(4)}} \text{translation by }(h,g)\,dh dg$.
\end{thm}

\begin{proof}
For $G=\GSp(4),\GL(2)$, denote by $\Meas\left(T^1_G(\bZ_p),V_{G,\ord}\right)$ the set of $p$-adic measures on $T^1_G(\bZ_p)$ valued in $V_{G,\ord}$. The group $T^1_G(\bZ_p)$ acts on it in two ways: via its action by translation on $T^1_G(\bZ_p)$ and its action on $V_{G,\ord}$. Let
\[
\Meas\left(T^1_G(\bZ_p),V_{G,\ord}\right)^\natural\subset \Meas\left(T^1_G(\bZ_p),V_{G,\ord}\right)
\]
be the subset on which the two action of $T^1_G(\bZ_p)$ agree and make it a $\tilde{\Lambda}_G$-module. Unfolding the definition of $\cM_{G,\ord}$ gives a natural map from $ \Meas\bigl(T^1_G(\bZ_p),V_{G,\ord}\bigr)^\natural$ to $\cM_{G,\ord}$ such that for each $p$-adic weight $\underline{\tau}\in \Hom_{\cont} \bigl(T^1_G(\bZ_p),\bar{\bQ}^\times_p \bigr)$, we~have the commutative diagram
\begin{equation}
\begin{tikzcd}
\Meas\left(T^1_G(\bZ_p),V_{G,\ord}\right)^\natural \arrow[r] \arrow[d,"\ts\text{evaluate at $\utau$}"']& \cM_{G,\ord} \arrow[d]\\
V_{G,\ord}[\utau] \arrow[r,"\ts\eqref{eq:embd}"] &\cM_{G,\ord}\otimes_{\tilde{\Lambda}_G}\tilde{\Lambda}_G/\cP_{\utau}
\end{tikzcd}
\end{equation}
(cf. \cite[\S6.1.4]{SLF}). Therefore, it suffices to show that inside the space,
\begin{equation}\label{eq:Meas-natural}
\Meas\Bigl((\bQ^\times\backslash\bA^\times_{\bQ,f}/U^p)\times T^1_{\GL(2)}\times T^1_{\GSp(4)},V_{\GL(2),\ord}\otimes_\cO V_{\GSp(4),\ord}\Bigr)^\natural,
\end{equation}
\begin{minipage}{\textwidth}
there exists a $p$-adic measure $\mu_\cE$ such that
\begin{multline}\label{eq:mu'}
\mu_\cE(\kappa,\tau,(\tau_1,\tau_2))
=e^{\GL(2)}_\ord e^{\GSp(4)}_\ord \,{\rm Proj}_{K^p_{\GL(2)},K^p_{\GSp(4)}}\bigl(E^\Sieg_{\kappa,\tau,\tau_1,\tau_2} \big|_{\GL(2)\times\GSp(4)}\bigr)\\
\times \spe_{\tau,(\tau_1,\tau_2)}(\omega^{-1}_{\sC_1}\omega^{-1}_{\sC_2}\circ{\rm det}_{\GL(2)})
\end{multline}
\end{minipage}
\par\smallskip\noindent
for all classical $(\kappa,\tau,\tau_1,\tau_2)=((k,\chi),(l,\xi),(l_1,l_2,\xi_1,\xi_2))$. The existence of $\mu_\cE$ can be shown by $p$-adic interpolation of Fourier coefficients.

Take sufficiently small tame level group $K^{p,\prime}_{\GL(2)}\subset K^p_{\GL(2)},K^{p,\prime}_{\GSp(4)}\subset K^p_{\GSp(4)}$ such that $\left.E^\Sieg_{\kappa,\tau,\tau_1,\tau_2}\right|_{\GL(2)\times\GSp(4)}$ is fixed by $K^{p,\prime}_{\GL(2)},K^{p,\prime}_{\GSp(4)}$ for all classical $(\kappa,\tau,\tau_1,\tau_2)$. Denote by $V'_{\GL(2)},V'_{\GSp(4)}$ the space of $p$-adic forms of tame levels $K^{p,\prime}_{\GL(2)},K^{p,\prime}_{\GSp(4)}$. Thanks to the strong approximation theorem for $\Sp(4)$ and $\SL(2)$, we~can pick $\nu_i\in\bA^{\times,p}_{\bQ,f}$, $i=1,\dots,c_1$, and $\nu'_j\in\bA^{\times,p}_{\bQ,f}$, $j=1,\dots,c_2$, such that each connected component of the Shimura variety for $\GL(2)\times\GSp(4)$ of level $K^{p,\prime}_{\GL(2)}\GL(2)(\bZ_p)\times K^{p,\prime}_{\GSp(4)}\GSp(4,\bZ_p)$ contains a cusp corresponding to $\bigl(\begin{smallbmatrix}\bid_2\\&\nu_i\cdot\bid_2\end{smallbmatrix},\begin{smallbmatrix}1\\&\nu'_j\end{smallbmatrix} \bigr)$ for some $i,j$. Taking the $q$-expansions at these cusps gives an injection
\begin{equation}\label{eq:q-exp}
\varepsilon_{\qexp}:V'_{\GL(2)}\hat{\otimes}_{\cO} V'_{\GSp(4)}
\hto \cO\llb \bQ_{>0}\times \Sym_2(\bQ)_{> 0}\rrb^{\oplus c_1c_2},
\end{equation}
and the image of $\varepsilon_{\qexp}$ is closed in $\cO\llb\Sym_2(\bQ)_{> 0}\times \bQ_{> 0}\rrb^{\oplus c_1c_2}$ (for the $p$-adic topology). (The $q$-expansion maps are defined via evaluations at Mumford objects \hbox{\cite[Ch.\,III]{FC}}, cf. \cite[\S3.4]{HidaCon}, \cite[\S4.2]{EisDiff}. The injectivity and the closedness of the image follows from the irreducibility of Igusa towers \cite[Cor.\,8.17]{HidaPAF}.) We view $V'_{\GL(2)}\hat{\otimes}_{\cO} V'_{\GSp(4)}$ as a subspace of $\cO\llb\Sym_2(\bQ)_{>0}\times \bQ_{> 0}\rrb^{\oplus c_1c_2}$. The $\bU_p$-operators on $V'_{\GL(2)}\hat{\otimes}_{\cO} V'_{\GSp(4)}$ extend to operators on $\cO\llb\Sym_2(\bQ)_{> 0}\times \bQ_{> 0}\rrb^{\oplus c_1c_2}$. Writing elements in $\cO\llb\Sym_2(\bQ)_{> 0}\times \bQ_{> 0}\rrb^{\oplus c_1c_2}$ as $\sum_{i,j,\beta_1,\beta_2} a_{(i,j)}(\beta_1,\beta_2) \,q^{(\beta_1,\beta_2)}_{(i,j)}$, with summation over \hbox{$1\leq i\leq c_1$}, $1\leq j\leq c_2$, $\beta_1,\beta_2\in \Sym_2(\bQ)_{>0}\times \bQ_{> 0}$, the extension has the formula
\begin{multline*}
U^{\GL(2)}_{p,m_3,0}U^{\GSp(4)}_{p,m_1,m_2,0}\biggl(\sum_{i,j,\beta_1,\beta_2} a_{i,j}(\beta_1,\beta_2)\,q^{(\beta_1,\beta_2)}\biggr)\\
=\sum_{x\in\bZ/p^{m_1-m_2}\bZ}\;
\sum_{i,j,\beta_1,\beta_2} a_{(i,j)}\left(\begin{smallbmatrix}p^{m_1-m_2}&\\ x&1\end{smallbmatrix}p^{m_2}\beta_1\begin{smallbmatrix}p^{m_1-m_2}&x\\&1\end{smallbmatrix},p^{2m_3}\beta_2\right)\,q^{(\beta_1,\beta_2)}.
\end{multline*}

On the other hand, letting $\cN_{G}$ denote the space of classical nearly holomorphic forms on $G=\GSp(4),\GL(2)$ over $\cO$ of all weights, invariant under the right translation by $K^{p,\prime}_G U_G(\bZ_p)$, and vanishing along all $p$-adic cusps, then the unit root splitting \cite[Th.\,4.1]{TDwork} gives rise to a map
\[
\imath_{p\adic}: \cN_{\GL(2)}\otimes_\cO\cN_{\GSp(4)}\to V'_{\GL(2)}\hat{\otimes}_{\cO} V'_{\GSp(4)},
\]
which is actually an embedding by \cite[Prop.\,3.12.1]{NHF}. Moreover, for a nearly holomorphic form $F$ on $\GL(2)\times\GSp(4)$, we~have
\begin{multline*}
(\beta_1,\beta_2)\text{-coefficient of }(i,j)\text{-component of }\varepsilon_{\qexp}\left(\imath_{p\adic} \left(F\right)\right)\\
= F_{\beta_1,\beta_2}\left(\begin{smallbmatrix}\bid_2\\&\nu_i\cdot\bid_2\end{smallbmatrix},\begin{smallbmatrix}1\\&\nu'_j\end{smallbmatrix}\right)(Y=0).
\end{multline*}
\begin{minipage}{\textwidth}
It follows that
\begin{multline*}
\varepsilon_{\qexp}\left(\imath_{p\adic} \left(E^\Sieg_{\kappa,\tau,\tau_1,\tau_2}\big|_{\GL(2)\times\GSp(4)}\right)\right)\\
= \sum_{\substack{(i,j),(\beta_1,\beta_2)\\ \nu_i=\nu'_j}} \Biggl(\sum_{\substack{\bbeta\in\Her_3(\cK)_{>0}\\[.5ex]\frac{\bbeta+\bar{\bbeta}}{2}=\begin{smallbmatrix}\beta_1&\ast\\ \ast&\beta_2\end{smallbmatrix}}}
\hspace{-1em}A(\bbeta;\kappa,\tau,\tau_1,\tau_2) \Bigl(\begin{Smallbmatrix}\bid_3\\&\nu_{i}\bid_3\end{Smallbmatrix}^p_f \Bigr) C_{k,l,l_1,l_2}(\bbeta) \Biggr) q^{(\beta_1,\beta_2)}_{(i,j)}
\end{multline*}
\end{minipage}
\par\smallskip\noindent
with $A(\bbeta;\kappa,\tau,\tau_1,\tau_2),C_{k,l,l_1,l_2}(\bbeta)$ given by the formulas \eqref{eq:Abeta} and \eqref{eq:Cbeta}. From \eqref{eq:Abeta}, it is easy to see that for each $\bbeta$ and $i,j$ such that $\nu_i=\nu'_j$, there exists a $p$-adic measure
\[
\mu_{(i,j),\bbeta}\in \Meas \Bigl(\bQ^\times\backslash\bA^\times_{\bQ,f}/U^p\times T^1_{\GL(2)}(\bZ_p)\times T^1_{\GSp(4)}(\bZ_p),\cO \Bigr)
\]
such that for call classical tuples $(\kappa,\tau,\tau_1,\tau_2)$,
\[
\mu_{(i,j),\bbeta}(\kappa,\tau,\tau_1,\tau_2)=A(\bbeta;\kappa,\tau,\tau_1,\tau_2) \Bigl(\begin{Smallbmatrix}\bid_3\\&\nu_i\bid_3\end{Smallbmatrix}^p_f \Bigr).
\]
Define
\[
\resizebox{\displaywidth}{!}{$\displaystyle \mu_{\cE,\qexp}\in \Meas\Bigl(\bQ^\times\backslash\bA^\times_{\bQ,f}/U^p\times T^1_{\GL(2)}(\bZ_p)\times T^1_{\GSp(4)}(\bZ_p),\cO\llb\bQ_{>0}\times \Sym_2(\bQ)_{> 0}\rrb^{\oplus c_1c_2}\Bigr)$} \]
as
\[
\mu_{\cE,\qexp}=\sum_{\substack{(i,j),(\beta_1,\beta_2)\\ \nu_i=\nu'_j}} \Biggl(\sum_{\bbeta\in\Her_3(\cK)_{>0},\,\frac{\bbeta+\bar{\bbeta}}{2}=\begin{bsm}\beta_1&\ast\\ \ast&\beta_2\end{bsm}}
\mu_{(i,j),\bbeta}\Biggr) q^{(\beta_1,\beta_2)}_{(i,j)}.
\]
Observe that $C_{k,l,l_1,l_2}(\beta)\equiv 1\bmod p^m$ for all $\bbeta$ such that $\bbeta+\bar{\bbeta}\equiv 0\bmod p^m$. Thus, for all classical $(\kappa,\tau,\tau_1,\tau_2)$,
\begin{multline}\label{eq:qexpE}
\left(U^{\GL(2)}_{p,1,0} U^{\GSp(4)}_{p,2,1,0}\right)^m \mu_{\cE,\qexp}(\kappa,\tau,\tau_1,\tau_2) \\
\equiv \varepsilon_{\qexp}\left( \left(U^{\GL(2)}_{p,1,0} U^{\GSp(4)}_{p,2,1,0}\right)^m\imath_{p\adic} \left(E^\Sieg_{\kappa,\tau,\tau_1,\tau_2}\big|_{\GL(2)\times\GSp(4)}\right)\right) \mod p^m.
\end{multline}

At a classical point, because $E^\Sieg_{\kappa,\tau,\tau_1,\tau_2}\big|_{\GL(2)\times\GSp(4)}$ belongs to a finite-rank $\cO$\nobreakdash-mod\-ule of classical nearly holomorphic forms of a certain level on which $U^{\GL(2)}_{p,1,0} U^{\GSp(4)}_{p,2,1,0}$ acts, the limit
\[
\lim_{n\to\infty} \left(U^{\GL(2)}_{p,1,0} U^{\GSp(4)}_{p,2,1,0}\right)^{n!} E^\Sieg_{\kappa,\tau,\tau_1,\tau_2}\big|_{\GL(2)\times\GSp(4)}
\]
exists. It follows that when $m=n!$ with $n\to\infty$, the limit of the right hand side of \eqref{eq:qexpE} exists. Combining this with the Zariski density of the classical points inside the weight space $\Hom_{\cont}\bigl((\bQ^\times\backslash\bA^\times_{\bQ,f}/U^p)\times T^1_{\GL(2)}\times T^1_{\GSp(4)},\ol{\bQ}^\times_p\bigr)$, we~deduce from \eqref{eq:qexpE} that the limit
\[
\lim_{n\to \infty}\left(U^{\GL(2)}_{p,1,0} U^{\GSp(4)}_{p,2,1,0}\right)^{n!} \mu_{\cE,\qexp}
\]
\begin{minipage}{\textwidth}
exists and interpolates $\varepsilon_{\qexp} \bigl(e^{\GL(2)}_\ord e^{\GSp(4)}_\ord\imath_{p\adic} \bigl(E^\Sieg_{\kappa,\tau,\tau_1,\tau_2}\big|_{\GL(2)\times\GSp(4)}\bigr)\bigr)$ at all classical $(\kappa,\tau,\tau_1,\tau_2)$. Denote this limit by
\begin{multline*}
e^{\GL(2)}_\ord e^{\GSp(4)}_\ord\left(\mu_{\cE,\qexp}\right)\\
\in\Meas\Bigl(\bQ^\times\backslash\bA^\times_{\bQ,f}/U^p\times T^1_{\GL(2)}(\bZ_p)\times T^1_{\GSp(4)}(\bZ_p),\cO\llb\bQ_{>0}\times \Sym_2(\bQ)_{> 0}\rrb^{\oplus c_1c_2}\Bigr).
\end{multline*}
\end{minipage}
Since the classical points are dense in the weight space and the image of \eqref{eq:q-exp} is closed, this limit must come from the $q$-expansion of a $p$-adic measure valued in $p$-adic forms, \ie there exists
\[
\mu'_\cE\in \Meas\Bigl((\bQ^\times\backslash\bA^\times_{\bQ,f}/U^p)\times T^1_{\GL(2)}\times T^1_{\GSp(4)},V'_{\GL(2),\ord}\otimes_\cO V'_{\GSp(4),\ord}\Bigr),
\]
such that
\[
\varepsilon_{\qexp}(\mu'_\cE)=e^{\GL(2)}_\ord e^{\GSp(4)}_\ord(\mu_{\cE,\qexp}).
\]
Moreover, by knowing that the evaluation of $\mu'_\cE$ at all classical points $(\kappa,\tau,\tau_1,\tau_2)$ has weight $\tau,(\tau_1,\tau_2)$, we~deduce that the two natural actions of $T^1_{\GL(2)}\times T^1_{\GSp(4)}$ on $\mu'_\cE$ agree with each other, and
\[
\mu'_\cE\in \Meas\Bigl((\bQ^\times\backslash\bA^\times_{\bQ,f}/U^p)\times T^1_{\GL(2)}\times T^1_{\GSp(4)},V'_{\GL(2),\ord}\otimes_\cO V'_{\GSp(4),\ord}\Bigr)^\natural.
\]
Then $\mu_\cE={\rm Proj}_{K^p_{\GL(2)},K^p_{\GSp(4)}}(\mu'_\cE)\cdot \omega^{-1}_{\sC_1}\omega^{-1}_{\sC_2}\circ{\rm det}_{\GL(2)}$ is the desired measure satisfying \eqref{eq:mu'}.
\end{proof}

\Subsection{The four-variable \texorpdfstring{$p$-adic \protect\Lfunction}{pL} and its interpolation formula I}

\subsubsection{Hida families and idempotent operators}

Let $F_{\sC_1},F_{\sC_2}$ be the function fields of the irreducible components $\sC_1,\sC_2$ fixed in Section~\ref{sec:setup}. Then the maps $\tilde{\Lambda}_{\GL(2)}\to F_{\sC_1}$, $\tilde{\Lambda}_{\GSp(4)}\to F_{\sC_2}$ factors through projections $\tilde{\Lambda}_{\GL(2)}\to\Lambda_{\GL(2)}$, $\tilde{\Lambda}_{\GSp(4)}\to\Lambda_{\GSp(4)}$ induced by characters of $T^1_{\GL(2)}(\bZ/p\bZ)$ and $T^1_{\GSp(4)}(\bZ/p\bZ)$. We~view $F_{\sC_1}$ as an algebra over $\Lambda_{\GL(2)}$ and $F_{\sC_2}$ as an $\Lambda_{\GSp(4)}$-algebra through these factorizations.

Denote by $\bI_{\sC_1},\bI_{\sC_2}$ the integral closures of $\Lambda_{\GL(2)},\Lambda_{\GSp(4)}$ inside $F_{\sC_1},F_{\sC_2}$. The universal ordinary Hecke algebras $\bT_{\GL(2),\ord},\bT_{\GSp(4),\ord}$ are known to be reduced. Therefore, we~have
\begin{align*}
\bT_{\GL(2),\ord}\otimes F_{\sC_1} = F_{\sC_1}\oplus R_{\sC_1},
\qquad\bT_{\GSp(4),\ord}\otimes F_{\sC_2} =F_{\sC_2}\oplus R_{\sC_2}
\end{align*}
as $F_{\sC_1}$-algebras and $F_{\sC_2}$-algebras such that the projection onto the first factor agrees with the natural maps $\bT_{\GL(2),\ord}\to \bI_{\sC_1}$, $\bT_{\GSp(4),\ord}\to \bI_{\sC_2}$. Let
\begin{align}
\label{eq:HeckeProj}
\mathds{1}_{\sC_1}\in \bT_{\GL(2),\ord}\otimes F_{\sC_1},
\qquad
\mathds{1}_{\sC_2}\in \bT_{\GSp(4),\ord}\otimes F_{\sC_2}
\end{align}
be the idempotent associated to the first factor in the above decomposition.

\subsubsection{The modified Euler factors at \texorpdfstring{$p$ and $\infty$}{pinf}}\label{sec:modEuler}
We let
\begin{align*}
\lambda_{\GL(2),0,1}\in\bI^\times_{\sC_1},
\qquad\lambda_{\GSp(4),0,0,1}, \lambda_{\GSp(4),1,0,0}\in \bI^\times_{\sC_2}
\end{align*}
denote the eigenvalues of the $\bU_p$-operators associated to $\begin{bsm}p\\&1\end{bsm},\begin{bsm}p\\&p\\&&1\\&&&1\end{bsm}, \begin{bsm}p\\&1\\&&p^{-1}\\&&&1\end{bsm}$.

Given a point $x\in \sC_1(\bar{\bQ}_p)\times\sC_2(\bar{\bQ}_p)$ where the weight projection map \eqref{eq:wtproj} is \'{e}tale and the image of $x$ is an arithmetic tuple $(\tau,\tau_1,\tau_2)=(l,\xi,l_1,\xi_1,l_2,\xi_2)$, we~let
\begin{align*}
\sS_{\GL(2),x}\qquad\text{(\resp $\sS_{\GSp(4),x}$)}
\end{align*}
be an orthogonal basis (with respect to the modified Petersson inner product defined in \eqref{eq:Pet2}, \eqref{eq:Pet}) of the space spanned by ordinary cuspidal holomorphic forms on $\GL(2)$ of weight $l$, tame level $K^p_{\GL(2)}$ (\resp $\GSp(4)$ of weight $(l_1,l_2)$, tame level $K^p_{\GSp(4)}$) and nebentypus at $p$ given by \eqref{eq:pneben}, belonging to the Hecke eigenspace parameterized by $x$. Let $\pi_x$ be the unitary irreducible automorphic representation of $\GL(2,\bA_\bQ)$ generated by forms $\sS_{\GL(2),x}$ twisted by a real power of $|\det|$, and $\Pi_x$ be a unitary irreducible automorphic representation of $\GSp(4,\bA_\bQ)$ inside the representation generated by forms in $\sS_{\GSp(4),x}$ twisted by a real power of $|\det|$. (There can be more than one choices of $\Pi_x$, but the partial \Lfunction and modified Euler factors at $p,\infty$ do not depend on the choice of $\Pi_x$.) Let $L^S(s,\Pi_x\times\pi_x\times\chi)$ be the degree $8$ partial L-function, and
\begin{multline} \label{eq:Ep}
E_p(s,\Pi_x\times\pi_x\times\chi)\\
=\gamma_p(s,\chi_p\eta_{x,1}\eta'_{x,1})^{-1} \gamma_p(s,\chi_p\eta_{x,1}\eta^{\prime}_{x,2})^{-1}
\gamma_p(s,\pi_{x,p}\times\chi_p\eta'_{x,3})^{-1},
\end{multline}
where the characters $\eta_{x,1},\eta_{x,2},\eta'_{x,1},\eta'_{x,2},\eta'_{x,3}$ are:
\begin{equation}\label{eq:etas}
\begin{aligned}
\eta_{x,1}(a)&=\xi(a|a|_p) \Bigl(p^{-\psfrac{l-1}{2}}\lambda_{\GL(2),0,1}(x)\Bigr)^{\val_p(a)},\\
\eta_{x,2}(a)&=\Bigl(p^{\psfrac{l-1}{2}}(\omega_{\sC_1,p}(p)\lambda^{-1}_{\GL(2),0,1})(x)\Bigr)^{\val_p(a)},\\
\eta'_{x,1}(a)&=\xi_1(a|a|_p)\Bigl(p^{\psfrac{-l_1+l_2-1}{2}} \omega_{\sC_2,p}(p)\lambda_{\GSp(4),1,0}\lambda_{\GSp(4),0,1}^{-1}(x)\Bigr)^{\val_p(a)},\\
\eta'_{x,2}(a)&=\xi_2(a|a|_p)\Bigl(p^{\psfrac{l_1-l_2+1}{2}} (\lambda_{\GSp(4),0,1}\lambda_{\GSp(4),1,0}^{-1})(x)\Bigr)^{\val_p(a)},\\
\eta'_{x,3}(a)&=\xi_1\xi_2(a|a|_p)\Bigl(p^{-\psfrac{l_1+l_2-3}{2}}\lambda_{\GSp(4),0,1}(x)\Bigr)^{\val_p(a)},
\end{aligned}
\end{equation}
and
\begin{multline}\label{eq:Einf}
E_\infty(s,\Pi_x\times\pi_x\times\chi)\\
=e^{-(4s+l_1+l_2+l)\cdot \sfrac{\pi i}{2}}\,\Gamma_\bC\Bigl(s+\frac{l_1+l_2+l}{2}-2\Bigr)\Gamma_\bC\Bigl(s+\frac{l_1+l_2-l}{2}-1\Bigr)\\
\times\Gamma_\bC\Bigl(s+\frac{-l_1+l_2+l}{2}-1\Bigr)\Gamma_\bC\Bigl(s+\frac{l_1-l_2+l}{2}\Bigr),
\end{multline}
where $\Gamma_\bC(s)=2(2\pi)^{-s}\Gamma(s)$.

The factors $E_p,E_\infty$ are obtained by unfolding the definitions in \cite{CoPerrin,Coates}. In \cite[p.\,163]{Coates}, for a homogeneous motive $M$ over $\bQ$ with coefficients in $K$ and $\sigma$ a fixed embedding $K\hto\bC$, the modified Euler factor at $p$ for its \Lfunction $L(\sigma,M,s)$ is defined as $\prod\nolimits_{U} E_p(\sigma,U,\rho,s)$, where $U$ runs over irreducible components of the semi-simplification of the Weil-Deligne representation at $p$ associated to the $\lambda$-adic realization of $M$ with $\lambda$ a prime in $K$ coprime to $p$, and $\rho$ is a choice of $\pm i$, for which we choose $\rho=i$. Interpreting this definition in our automorphic context here with $\Pi_x$ (\resp $\pi_x$) a subquotient of the induced representation associated to
\begin{align*}\arraycolsep4.5pt
\begin{bmatrix} a_1&\ast&\ast&\ast\\& a_2&\ast&\ast\\ &&\delta a_2\\&&\ast&\delta a_1\end{bmatrix}&\mto \eta'_{x,1}(a_1)\,\eta'_{x,2}(a_2)\,\eta'_{x,3}(\delta)\\
\arraycolsep4.5pt
\Bigl(\text{\resp }\begin{bmatrix}a&\ast\\&d\end{bmatrix}&\mto \eta_{x,1}(d)\,\eta_{x,2}(a)\Bigr),
\end{align*}
the irreducible component $U$ corresponds to $\eta=\eta_{x,i}\eta'_{x,j},\eta_{x,i}\eta'_{x,1}\eta'_{x,2}\eta^{\prime,-1}_{x,3}$, $i=1,2$, $j=1,2,3$. Denoting by $\Phi$ an element in $G_{\bQ_p}$ whose image in $G_{\bQ_p}/I_p$ is the geometric Frobenius, the factor $E_p(\sigma,U,\rho,s)$ is defined to be $1$ or the inverse of the $\gamma$-factor of~$U$ depending on whether the inverse roots of $\det(1-\Phi.X|U)$ have $p$-adic valuations larger or smaller than $-\sfrac{1}{2}$, namely in our context here whether $\val_p(\eta(p)p^{-s})$ is larger or smaller than $-\sfrac{1}{2}$. Note that the $\lambda$'s in \eqref{eq:etas} are all $p$-adic units. Thus, when
\[
-\frac{\min\{-l_1+l_2+l,\,l_1+l_2-l\}}{2}+2 \leq s\leq \frac{\min\{-l_1+l_2+l,\,l_1+l_2-l\}}{2}-1,
\]
among all the $\eta$'s, the ones with $\val_p(\eta(p)p^{-s})<-\sfrac{1}{2}$ are exactly
\[
\eta_{x,1}\eta'_{x,1},\,\eta_{x,1}\eta'_{x,2},\,\eta_{x,1}\eta'_{x,3},\,\eta_{x,2}\eta'_{x,3}.
\]
Therefore, the modified Euler factor at $p$ is given as \eqref{eq:Ep}. In \cite[(38) on p.\,157]{Coates}, the modified Euler factor at $\infty$ for $M$ is defined as $\prod_U E_\infty(U,\rho,s)$ where $U$ runs over the direct summands of the Hodge decomposition of $H_B(M)\otimes_{K,\sigma}\bC$. In our case here, we~only have $U$'s of type (a) in \loccit where $E_\infty(U,\rho,s)$ can be obtained from $L_\infty(U,s)$ by replacing $\Gamma_\bC(s)$ with $\Gamma_{\bC,\rho}(s)=\rho^{-s}\Gamma_\bC(s)$. It follows that \hbox{$E_\infty(s,\Pi_x\times\pi_x\times\chi)$} equals the product of\enlargethispage{-2\baselineskip}
\begin{multline*}
L_\infty(s,\Pi_x\times\pi_x\times\chi)=\Gamma_\bC\Bigl(s+\frac{l_1+l_2+l}{2}-2\Bigr)\Gamma_\bC\Bigl(s+\frac{l_1+l_2-l}{2}-1\Bigr)\\
\times\Gamma_\bC\Bigl(s+\frac{-l_1+l_2+l}{2}-1\Bigr)\Gamma_\bC\Bigl(s+\frac{l_1-l_2+l}{2}\Bigr),
\end{multline*}
and $\rho^{-(s+\psfrac{l_1+l_2+l}{2}-2)-(s+\psfrac{l_1+l_2-l}{2}-1)-(s+\psfrac{-l_1+l_2+l}{2}-1)-(s+\psfrac{l_1-l_2+l}{2})}$.

\subsubsection{The modified Petersson inner product and Bessel period}\label{sec:modper}
For ordinary holomorphic automorphic forms $f_1,f_2$ on $\GL(2)$ and $\varphi_1,\varphi_2$ on $\GSp(4)$, we~define the modified Petersson inner product $\bfP(f_1,f_2),\bfP(\varphi_1,\varphi_2)$ and the modified Bessel \hbox{period}~$\bes^\dagger_{\bS,\Lambda}(\varphi_1)$ as follows:
\begin{align}\label{eq:Pet2}
&\bfP(f_1,f_2)\\
&\hspace*{3mm}=\lambda_p(f_1)^{-m}\int_{[\GL(2)]} f_1(g)\,f_2\bigl(g\begin{smallbmatrix}&1\\1\end{smallbmatrix}_{p\infty} \begin{smallbmatrix}p^{m}\\&p^{-m}\end{smallbmatrix}_p\bigr)\cdot\omega_\pi(\det g)^{-1}\,dg,\notag
\\
&\bfP(\varphi_1,\varphi_2)=\lambda_{p,1}(\varphi_1)^{-m_1}\lambda_{p,2}(\varphi_1)^{-m_2} \label{eq:Pet} \\
&\hspace*{3mm}\times\int_{[\GSp(4)]} \varphi_1(g)\,\varphi_2\Biggl(g\begin{bsm}&\bid_2\\\bid_2\end{bsm}_{p\infty}\begin{bsm}p^{m_1}\\&p^{m_2}\\&&p^{-m_1}\\&&&p^{-m_2}\end{bsm}_p\Biggr)\cdot \omega_\Pi(\nu_g)^{-1} \,dg,\notag\\
&\bes^\dagger_{\bS,\Lambda}(\varphi_1)=
\lambda_{p,1}(\varphi_1)^{-m_1}\lambda_{p,2}(\varphi_1)^{-m_2}
\cdot \bes_{\bS,\Lambda}\Biggl(\begin{bsm}p^{m_1}\\&p^{m_2}\\&&p^{-m_1}\\&&&p^{-m_2}\end{bsm}_p\varphi_1\Biggr),\label{eq:bes-dagger}
\end{align}
with $m_1\gg m_2\gg 0$, $m\gg 0$, where $\lambda_p(f_1),\lambda_{p,1}(\varphi_1),\lambda_{p,2}(\varphi_1)$ are defined by
\begin{align*}
\lambda_p(f_1)^m f_1&=\int_{U_{\GL(2)}(\bZ_p)} \pi_p\left(u\begin{smallbmatrix}p^m\\&p^{-m}\end{smallbmatrix}\right)f_1\,du,\\
\lambda_{p,1}(\varphi_1)^{m_1}\lambda_{p,2}(\varphi_1)^{m_2} \varphi_1&=\int_{U_{\GSp(4)}(\bZ_p)} \Pi_p\left(u\begin{bsm}p^{m_1}\\&p^{m_2}\\&&p^{-m_1}\\&&&p^{-m_2}\end{bsm}\right)\varphi_1\,du,
\end{align*}
and $\bes_{\bS,\Lambda}$ is the Bessel period as defined in \cite[\S2.2.1]{pFL}. (One can check that the right hand sides of \eqref{eq:Pet2}, \eqref{eq:Pet}, \eqref{eq:bes-dagger} do not depend on $m_1,m_2,m$ as long as $m_1-m_2,m_2,m$ are sufficiently large. See \cite[Prop.\,2.7.1]{pFL} for a proof of this for \eqref{eq:bes-dagger}.)

\subsubsection{The four-variable \texorpdfstring{$p$-adic \protect\Lfunction}{pL}}

With the various objects defined in Sections~\ref{sec:modEuler}, \ref{sec:modper}, we~apply the idempotent in \eqref{eq:HeckeProj} to the $p$-adic measure $\mu_\cE$ constructed in Theorem~\ref{thm:EisF} to obtain the $p$-adic measure for $\sC_1,\sC_2$. Before stating the theorem, we~introduce some more notation.

With $\bS=\begin{smallbmatrix}\bba&\sfrac{\bbb}{2}\\ \sfrac{\bbb}{2}&\bbc\end{smallbmatrix}\in\Sym_2(\bQ)_{>0}
$ and $\bLam:\cK^\times\backslash\bA^\times_{\cK,f}\to \Lambda^\times_{\GSp(4)}$ as fixed in Section~\ref{sec:setup}, for $f\in\sS_{\GL(2),x},\varphi\in\sS_{\GSp(4),x}$, we~let $f_\bbc$ denote the Fourier coefficients of~$f$ indexed by $\bbc$, and $\bes^\dagger_{\bS,\Lambda}\left(\varphi\right)$ denote the modified Bessel period in \eqref{eq:bes-dagger} with $\Lambda$ the classical Hecke character corresponding to the specialization of $\bLam$ at $(l_1,l_2,\xi_1,\xi_2)$. We~denote by $(r_{\sLambda,1},r_{\sLambda,2})$ the $\infty$-type of $\Lambda$.

\begin{thm}\label{thm:pFL1}
Given the data in Section~\ref{sec:setup}, there exists
\[
\mu^S_{\sC_1,\sC_2}\!\in\!
\Meas\Bigl(\bQ^\times\backslash\bA^\times_{\bQ,f}/U^p,\cM_{\GL(2),\ord}\wh{\otimes}_{\cO} \cM_{\GSp(4),\ord}\Bigr)
\otimes_{\tilde{\Lambda}_{\GL(2)}\hat{\otimes}_\cO\tilde{\Lambda}_{\GSp(4)}}\!(F_{\sC_1}\hat{\otimes}_\cO F_{\sC_2})
\]
satisfying the interpolation properties: Suppose that $x\in \sC_1(\bar{\bQ}_p)\times\sC_2(\bar{\bQ}_p)$ is a point at which the weight projection map \eqref{eq:wtproj} is \'{e}tale. Then $ \mu^S_{\sC_1,\sC_2}$ has no poles along $x$. Let $\tau\in\Hom_\cont\bigl(T^1_{\GL(2)}(\bZ_p),\bar{\bQ}_p\bigr)$, $(\tau_1,\tau_2)\in \Hom_\cont\bigl(T^1_{\GSp(4)}(\bZ_p),\bar{\bQ}_p\bigr)$ be the projection of $x$ to the weight space.
\par\noindent
\begin{minipage}{\textwidth}
For a character $\kappa\in \Hom_\cont\bigl(\bQ^\times\backslash\bA^\times_{\bQ,f}/U^p,\bar{\bQ}_p\bigr)$ such that $(\kappa,\tau,\tau_1,\tau_2)$ is classical (as~defined at the end of Section~\ref{sec:pchar}),
\begin{multline}\label{eq:partial-int}
\mu^S_{\sC_1,\sC_2}(\kappa,x)\\
=\sum_{f\in \sS_{\GL(2),x}} \frac{f_{\bbc}f}{\bfP(f,f)}\sum_{\varphi\in \sS_{\GSp(4)},x} \frac{\bes^\dagger_{\bS,\Lambda}(\varphi)\varphi}{\bfP(\varphi,\varphi)}
\cdot i^{r_{\sLambda,1}-r_{\sLambda,2}} I_\infty(k,\cD_{l_1,l_2},\cD_l,\Lambda_\infty)\\
\times E_p\Bigl(k+\frac{l+l_1+l_2}{2},\Pi_x\times\pi_x\times \chi\Bigr)\cdot L^S\Bigl(k+\frac{l+l_1+l_2}{2},\Pi_x\times\pi_x\times \chi\Bigr),
\end{multline}
\end{minipage}
\par\smallskip\noindent
with $I_\infty(k,\cD_{l_1,l_2},\cD_l,\Lambda_\infty)$ the Archimedean zeta integral given in \cite[(4.2.5)]{pFL} with $k+\sfrac{\epsilon}{2}$, $t_k$ in \loccit equal to $k+\psfrac{l+l_1+l_2}{2}$, $|2k+l+l_1+l_2-1|+3$ here and $\cD_{l_1,l_2},\cD_l$ holomorphic discrete series of $\GSp(4),\GL(2)$ of weights $(l_1,l_2),l$. If~further assuming that $l=l_1=l_2$, then
\begin{multline*}
\mu^S_{\sC_1,\sC_2}(\kappa,x)
=\bbc\sqrt{\det\bS}\,2^{-3l+2}i^l\sum_{f\in \sS_{\GL(2),x}} \frac{f_{\bbc}f}{\bfP(f,f)}\sum_{\varphi\in \sS_{\GSp(4)},x} \frac{\bes^\dagger_{\bS,\Lambda}\left(\varphi\right)\varphi}{\bfP(\varphi,\varphi)} \\
\times E_\infty \Bigl(k+\frac{3l}{2},\Pi_x\times\pi_x\times \chi\Bigr) E_p\Bigl(k+\frac{3l}{2},\Pi_x\times\pi_x\times \chi\Bigr) L^S\Bigl(k+\frac{3l}{2},\Pi_x\times\pi_x\times \chi\Bigr).
\end{multline*}
(See Sections~\ref{sec:modEuler}, \ref{sec:modper} for the definitions of the terms appearing in the above inter\-polation formulas.)
\end{thm}

\begin{proof}
We first examine the evaluations of
\begin{multline*}
(\mathds{1}_{\sC_1}\otimes \mathds{1}_{\sC_2})\cdot \mu_{\cE} \in \Meas\left(\bQ^\times\backslash\bA^\times_{\bQ,f}/U^p,\cM_{\GL(2),\ord}\wh{\otimes}_{\cO} \cM_{\GSp(4),\ord}\right)\\
\otimes_{\tilde{\Lambda}_{\GL(2)}\hat{\otimes}_\cO\tilde{\Lambda}_{\GSp(4)}} (F_{\sC_1}\hat{\otimes}_\cO F_{\sC_2}).
\end{multline*}

By the construction of $\mu_\cE$, we~know that at $(\kappa,x)$ with $(\kappa,\tau,\tau_1,\tau_2)$ classical,
\[
\left((\mathds{1}_{\sC_1}\otimes \mathds{1}_{\sC_2})\cdot \mu_{\cE}\right)(\kappa,x)
=\hspace*{-3mm}\sum_{f\in \sS_{\GL(2),x}}\sum_{\varphi\in \sS_{\GSp(4),x}}\hspace*{-5mm}
\frac{\left.\bfP\left(E^\Sieg_{\kappa,\tau,\tau_1,\tau_2}\right|_{\GL(2)\times\GSp(4)},f\otimes\varphi\right)}{\bfP(f,f)\bfP(\varphi,\varphi)}
f\otimes\varphi.
\]
(Here we use that the definition of $\bfP(-,-)$ plus the tame level and ordinarity of $f,\varphi$ implies that applying $\bfP(-,f\otimes \varphi)$ to $\left.E^\Sieg_{\kappa,\tau,\tau_1,\tau_2}\right|_{\GL(2)\times\GSp(4)}$ and
\[
\mu_\cE(\kappa,x)=e^{\GL(2)}_\ord e^{\GSp(4)}_\ord \,{\rm Proj}_{K^p_{\GL(2)},K^p_{\GSp(4)}}\Bigl(\left.E^\Sieg_{\kappa,\tau,\tau_1,\tau_2}\right|_{\GL(2)\times\GSp(4)}\Bigr)
\]
produces the same value.) The computation in the proof of \cite[Th.\,4.2.1]{pFL} (or more precisely the formula for $I_p(s)$ and $C_{k,\chi,\Pi,\pi}(s)$ in \loccit gives
\begin{multline*}
\bfP\bigl(E^\Sieg_{\kappa,\tau,\tau_1,\tau_2},f\otimes\varphi\big|_{\GL(2)\times\GSp(4)}\bigr)
=C(k,\chi,x)\cdot f_\bbc \bes^\dagger_{\bS,\Lambda}(\varphi)
\cdot I_\infty(k,\Pi_{x,\infty},\pi_{x,\infty},\Lambda_\infty)\\
\times E_p\Bigl(k+\frac{l+l_1+l_2}{2},\Pi_x\times\pi_x\times \chi\Bigr)\cdot L^S\Bigl(k+\frac{l+l_1+l_2}{2},\Pi_x\times\pi_x\times, \chi\Bigr)
\end{multline*}
\begin{minipage}{\textwidth}
and
\begin{multline}\label{eq:factor-C}
C(k,\chi,x)\\
=\vol_S \, \frac{1-\left(\sfrac{\cK}{p}\right)p^{-1}}{1+p^{-1}}\,\biggl|\frac{\alphaS-\alphabS}{\sqrt{\disc(\cK/\bQ)}}\biggr|_p\cdot |\bbc|^2_p |(\alphaS-\alphabS)^2|^2_p\cdot \chi_p(-1)(-1)^k\\
\shoveleft{\hspace*{1mm}\times(\chi^{-1}_p|\cdot|^{-k}_p)(\bbc) \bbc^{-k} \cdot \xi^{-1}\xi^{-1}_1\xi^{-1}_2(\bbc|\bbc|_p) (\bbc|\bbc|_p)^{-l-l_1-l_2}\cdot \lambda_{\GL(2)}\lambda_{\GSp(4)}(x)^{-\val_p(\bbc)}} \\
\shoveleft{\hspace*{.2mm}\times(\chi^{-1}_p|\cdot|^{-k}_p)(-(\alphaS-\alphabS)^2)(-(\alphaS-\alphabS)^2)^{-k}} \\
\shoveleft{\hspace*{.2mm}\times \xi^{-1}(-(\alphaS-\alphabS)^2|\alphaS-\alphabS|^2_p)(-(\alphaS-\alphabS)^2|\alphaS-\alphabS|^2_p)^{-l}}\\
\shoveleft{\hspace*{.2mm}\times\Lambda^{-c}_p(-(\alphaS-\alphabS))(\alphaS-\alphabS)^{-r_{\sLambda_1}}(-\alphaS+\alphabS)^{-r_{\sLambda_2}}}\\
\shoveleft{\hspace*{.2mm}\times \lambda_{\GL(2)}(x)^{-\val_p((\alphaS-\alphabS)^2)}\cdot i^{r_{\sLambda,1}-r_{\sLambda,2}},\hfill\mbox{}}
\end{multline}
\end{minipage}
\par\smallskip\noindent
where $\vol_S$ is a nonzero constant independent of $k,\chi,x$ (and can be expressed in terms of the volumes of some open compact subgroups of $\GL_2(\bQ_v),\GSp(4,\bQ_v)$ at finite $v\in S$), $\lambda_{\GL(2)}\in \bI^\times_{\sC_1},\lambda_{\GSp(4)}\in\bI^\times_{\sC_2}$ denote the eigenvalues of the $\bU_p$-operators corresponding to $\begin{smallbmatrix}p\\&1\end{smallbmatrix},\begin{Smallbmatrix}p\cdot\bid_2\\&\bid_2\end{Smallbmatrix}$ along $\sC_1,\sC_2$. (To plug in the formulas in \cite{pFL}, $s$, $k+\sfrac{\epsilon}{2}$, $\Lambda,\eta_{\pi_p,1}(a)$, $\eta_{\sPi_p,3}(a)$ for $a\in\bQ_p$ in \loccit correspond to
\begin{gather*}
k+\frac{l+l_1+l_2-1}{2},\quad k+\frac{l+l_1+l_2}{2},\\
{\Lambda^{-c}|\cdot|^{\psfrac{l_1+l_2}{2}}_{\bA_\cK}},\; \xi^{-1}(a|a|_p)\cdot(p^{\psfrac{l-1}{2}}\lambda^{-1}_{\GL(2)}(x))^{\val_p(a)},\\\xi^{-1}_1\xi^{-1}_2(a|a|_p)\cdot (p^{\psfrac{l_1+l_2-3}{2}}\lambda^{-1}_{\GSp(4)}(x))^{\val_p(a)}
\end{gather*}
here. Also, note that the integral in \loccit is over $\GU(1,1)\times\GSp(4)$ with $f$ extended to $\GU(1,1)$ by $\Upsilon$ equals our integral over $\GL(2)\times\GSp(4)$ here because the central characters match.)

From the above formula for $C(k,\chi,x)$, it is easy to see that there exists
\[
\cC\in \left(\cO\llb \bQ^\times\backslash\bA^\times_{\bQ,f}/U^p\rrb \hat{\otimes}_{\cO} \bI_{\sC_1}\hat{\otimes}_\cO \bI_{\sC_2}\otimes_\cO F\right)^\times
\]
such that for all $(\kappa,x)$ with $(\kappa,\tau,\tau_1,\tau_2)$ classical,
\[
\cC(\kappa,x)=i^{-r_{\sLambda,1}+r_{\sLambda,2}}\cdot C(k,\chi,x).
\]
Let
\begin{align*}
\mu^S_{\sC_1,\sC_2}= \cC^{-1}\left((\mathds{1}_{\sC_1}\otimes \mathds{1}_{\sC_2})\cdot \mu_{\cE}\right).
\end{align*}
Then for $(\kappa,x)$ as in the statement of the theorem,
\begin{multline*}
\mu^S_{\sC_1,\sC_2}(\kappa,x)\\
=\sum_{f\in \sS_{\GL(2),x}}\sum_{\varphi\in \sS_{\GSp(4),x}}
\frac{f_\bbc \bes^\dagger_{\bS,\Lambda}(\varphi)}{\bfP(f,f)\bfP(\varphi,\varphi)}
\cdot i^{r_{\sLambda,1}-r_{\sLambda,2}}I_\infty(k,\cD_{l_1,l_2},\cD_l,\Lambda_\infty)\\
\times E_p\Bigl(k+\frac{l+l_1+l_2}{2},\Pi_x\times\pi_x\times \chi\Bigr)\cdot L^S\Bigl(k+\frac{l+l_1+l_2}{2},\Pi_x\times\pi_x\times \chi\Bigr)
\cdot f\otimes\varphi.
\end{multline*}
When $l_1=l_2=l$, the integral $I_\infty(k,\cD_{l_1,l_2},\cD_l,\Lambda_\infty)$ is evaluated in \cite[Prop.\,4.2]{pFL}. Like in \loccit, plugging the result into the right hand side of the above formula gives the complete interpolation formula in the case $l_1=l_2=l$.
\end{proof}

\section{Specialization to Hida families of Yoshida lifts}
The interpolation formula in Theorem~\ref{thm:pFL1} is incomplete in the sense that it contains an uncomputed Archimedean zeta integral $I_\infty(k,\cD_{l_1,l_2},\cD_l,\Lambda_\infty)$ when $l_1,l_2,l$ are not all equal. In order to calculate this integral, we~let $\sC_2=\theta(\sB,\sB')$, the Hecke eigensystem associated to the Yoshida lifts of Hida families $\sB,\sB'$ on $\GL(2)$, and compare $\mu^S_{\sC_1,\theta(\sB,\sB')}$ with the product of Rankin--Selberg $p$-adic \Lfunctions for $\sB,\sC_1$ and $\sC_1,\sB'$. The main difficulty in carrying out this strategy lies in comparing the periods, which boils down to comparing the (modified) Petersson norms of a pair of modular forms and that of their Yoshida lift. To circumvent the explicit calculation of the corresponding Rallis inner product formula in our specific case, we~make use of the standard $p$-adic \Lfunction for $\theta(\sB,\sB')$, which has the (modified) Petersson norms of Yoshida lifts as periods.

\subsection{Some previous results on \texorpdfstring{$p$-adic \protect\Lfunctions}{pL}}\label{sec:pL}

For simplicity, in this section and Section~\ref{sec:compare}, we~assume that $S$ contains some finite place $v\neq p$, and for all such $v\in S$,
\[
K^p_{\GL(2),v}=\bigl\{g\in\GL(2,\bZ_v):g\equiv \begin{smallbmatrix}\ast&\ast\\0&\ast\end{smallbmatrix}\bmod \varpi_v\bigr\}.
\]

We recall some previous results on constructions of Kubota--Leopold $p$-adic \Lfunctions, Rankin--Selberg $p$-adic \Lfunctions, and $p$-adic (degree 5) standard \Lfunctions for $\Sp(4)$. Complete interpolation properties have been established for these $p$-adic \Lfunctions, and they will be used to formulate the right hand side in the key identity~\eqref{eq:compare}. To simplify the writing of the interpolation properties, we~use the following convention: Given an automorphic representation $\sigma$ with $\sigma_\infty$ isomorphic to holomorphic discrete series and ordinary at $p$, we~let
\[
D^S(s,\sigma)=E_\infty(s,\sigma)\,E_p(s,\sigma)\,L^S(s,\sigma)
\]
with $E_\infty(s,\sigma),E_p(s,\sigma)$ the modified Euler factor at $\infty$ and $p$ for $p$-adic interpolation as defined in \cite{CoPerrin,Coates}.

\Subsubsection{Kubota--Leopoldt \texorpdfstring{$p$-adic \protect\Lfunction}{pL}}
\begin{thm}\label{thm:KLp}
Given a Dirichlet character $\phi$ unramified outside $S\setminus\{p\}$, there exists
\[
\cL^S_{\KL,\phi}\in \Meas\bigl(\bQ^\times\backslash\bA^\times_{\bQ,f}/\hat{\bZ}^{p,\times},\cO\bigr)
\]
such that for all arithmetic $\kappa=(k,\chi)\in \Hom_\cont\bigl(\bQ^\times\backslash\bA^\times_{\bQ,f}/\hat{\bZ}^{p,\times},\bar{\bQ}_p\bigr)$ with $\phi\chi(-1)=(-1)^k$ and $k\geq 1$ or $\phi\chi(-1)=(-1)^{k+1}$ and $k\leq 0$,
\begin{align*}
\cL^S_{\KL,\phi}(\kappa)=D^S(k,\phi\chi).
\end{align*}
\end{thm}
(Here, $S$ is assumed to contain finite places other than $p$. Hence, the imprimitive Kubota--Leopoldt $p$-adic \Lfunction does not have poles.)

\Subsubsection{Rankin--Selberg \texorpdfstring{$p$-adic \protect\Lfunction}{pL}}

\begin{thm}\label{thm:Rankin}
Let $\sB_1,\sB_2\subset \Spec(\bT_{\GL(2),\ord})$ be two geometrically irreducible components. We~assume that $\sB_1$ is {\it primitive}, \ie the newforms in the automorphic representations corresponding to classical points of $\sB_1$ has tame level equal to $K^p_{\GL(2)}$). Denoting by $F_{\sB_1},F_{\sB_2}$ the function fields of $\sB_1,\sB_2$, there exists
\[
\cL^S_{\sB_1,\sB_2}\in \Meas\left(\bQ^\times\backslash\bA^\times_\bA/U^p,\Lambda_{\GL(2)}\hat{\otimes}_\cO\Lambda_{\GL(2)}\right)\otimes_{\Lambda_{\GL(2)}\hat{\otimes}_\cO\Lambda_{\GL(2)}} (F_{\sB_1}\hat{\otimes}_\cO F_{\sB_2})
\]
satisfying the interpolation property: Suppose that $(x_1,x_2)\in \sB_1(\bar{\bQ}_p)\times\sB_2(\bar{\bQ}_p)$ is a classical point of weights $t_1> t_2\geq 2$. Then for an arithmetic character $\kappa=(k,\chi)\in \Hom_\cont\bigl(\bQ^\times\backslash\bA^\times_{\bQ,f}/U^p,\bar{\bQ}^\times_p\bigr)$ such that
\[
-t_1+1\leq k\leq -t_2
\]
\ie $s=k+\psfrac{t_1+t_2}{2}$ is a critical point for the Rankin--Selberg \Lfunction \hbox{$L(s,\sigma_{x_1}\times\sigma_{x_2})$}, we~have
\begin{align}
\label{eq:RS-int} \cL^S_{\sB_1,\sB_2}(\kappa,x)=\frac{D^S\left(k+\psfrac{t_1+t_2}{2},\sigma_{x_1}\times\sigma_{x_2}\times\chi\right)}{ (-2i)^{t_1+1}\bfP(f_{x_1},f_{x_1})},
\end{align}
where $\sigma_{x_j}$, $j=1,2$, is the (unique) unitary automorphic representation of $\GL(2)$ (with unitary central character) giving rise to the Hecke eigensystem parameterized by $x_j$ (up to a twist by a real power of $|\det|$), and $f_{x_1}\in\sigma_{x_1}$ is the normalized eigenform for the Hecke eigensystem parameterized by $x_1$. (The modified Petersson inner product $\bfP(f_{x_1},f_{x_1})$ is defined as in \eqref{eq:Pet2}.)
\end{thm}

This theorem is proved in \cite[Th.\,5.1d]{Hi88}, see also \cite[Th.\,A]{Chen-Hsieh}. Our $\cL^S_{\sB_1,\sB_2}$ are obtained from the $p$-adic \Lfunctions in \loccit by removing $L$-factors at $v\in S\setminus\{p,\infty\}$ which are $p$-adically interpolatable and by a change of variable. Note that the classical Petersson norm for the new form inside $\sigma_{x_1}$ is used in the interpolation formula in \loccit, while in the above theorem, the modified Petersson norm $\bfP(f_{x_1},f_{x_1})$ is used. The relation between the two Petersson norms are proved in \cite[Lem.\,3.6]{HsiehTriple}, via which one can see that the interpolation formula \eqref{eq:RS-int} follows from that in \cite{Hi88,Chen-Hsieh}. (We also note that there is a slight difference between our convention of nebentypus here and that in \loccit If $\sigma_j$ has central character $\omega_j$, then in our convention the $p$-nebentypus is $\omega_j|_{\bZ^\times_p}$ and in the conventions in \loccit, the nebentypus sends $q_v$ to $\omega_v(q_v)$, so essentially is $\omega^{-1}_j|_{\bZ^\times_p}$.)

\subsubsection{Standard \texorpdfstring{$p$-adic \protect\Lfunction}{pL} for Yoshida lifts}

Let $\sB,\sB'\subset \Spec(\bT_{\GL(2),\ord})$ be two geometrically irreducible components. Let $F_{\sB},F_{\sB'}$ be the function fields of $\sB,\sB'$ and $\bI_{\sB},\bI_{\sB'}$ be the integral closures of $\Lambda_{\GL(2)}$ in $F_{\sB},F_{\sB'}$. Denote by
\begin{align*}
\lambda_{\sB}:\bT_{\GL(2),\ord}\to \bI_{\sB},
\qquad\lambda_{\sB'}:\bT_{\GL(2),\ord}\to \bI_{\sB'}
\end{align*}
the corresponding Hecke eigensystems, and
\begin{align*}
\omega_{\sB},\omega_{\sB'}:\bQ^\times\backslash\bA^\times_{\bQ,f}\to \Lambda^\times_{\GL(2)}
\end{align*}
the central characters. We~fix a square root for each of them
\[
\omega^{\sfrac{1}{2}}_{\sB},\omega^{\sfrac{1}{2}}_{\sB'}:\bQ^\times\backslash\bA^\times_{\bQ,f}\to \Lambda^\times_{\GL(2)}.
\]
(The equality proved in Proposition~\ref{prop:compare} below does not depend on the choice of the square roots.)

We have the group homomorphism
\begin{align*}
T^1_{\GL(2)}\times T^1_{\GL(2)}&\to T^1_{\GSp(4)},\\
\bigl(\diag(a_1,a^{-1}_1),\diag(a_2,a^{-1}_2)\bigr)&\mto \diag(a_1a_2,a_1a^{-1}_2,a^{-1}_1a^{-1}_2,a^{-1}_1a_2)
\end{align*}
which induces
\[
\wt{\Lambda}_{\GL(2)}\hat{\otimes}\wt{\Lambda}_{\GL(2)} \to \wt{\Lambda}_{\GSp(4)}.
\]
Let
\[
\bI_{\theta(\sB,\sB')}= \wt{\Lambda}_{\GSp(4)}\otimes_{\wt{\Lambda}_{\GL(2)}\hat{\otimes}\wt{\Lambda}_{\GL(2)}} (\bI_{\sB}\hat{\otimes}_{\cO}\bI_{\sB'}).
\]
It follows from the theory of theta lifts \cite{RobGlobalTheta} that if there exists a finite place $v\neq p$ such that the classical specializations of $\sB,\sB'$ are discrete series at $v$, then
for suitable $K^p_{\GSp(4)}$, there exists a geometrically irreducible component $\theta(\sB,\sB')\subset \Spec(\bT_{\GSp(4),\ord})$ with the $\wt{\Lambda}_{\GSp(4)}$-algebra homomorphism
\begin{align*}
\lambda_{\theta(\sB,\sB')}:\bT_{\GSp(4),\ord}\to \bI_{\theta(\sB,\sB')}
\end{align*}
such that the central character equals $\bQ^\times\backslash\bA^\times_f\To{\omega_{\sB}}\Lambda^\times_{\GL(2)}\to \Lambda^\times_{\GSp(4)}$ where the second map is induced by $T^1_{\GL(2)}\to T^1_{\GSp(4)}, \,\diag(a,a^{-1})\mto \diag(a,a,a^{-1},a^{-1})$, and
\begin{align*}
\lambda_{\theta(\sB,\sB')} \Bigl(\GSp(4,\bZ_v)\begin{Smallbmatrix}\varpi_v\\ &\varpi_v\\&&1\\&&&1\end{Smallbmatrix}\GSp(4,\bZ_v) \Bigr)
&=\lambda_{\sB}\bigl(\GL(2,\bZ_v)\begin{smallbmatrix}\varpi_v\\&1\end{smallbmatrix}\GL(2,\bZ_v)\bigr),\\
\lambda_{\theta(\sB,\sB')}\Bigl(\GSp(4,\bZ_v)\begin{Smallbmatrix}\varpi_v\\ &1\\&&1\\&&&\varpi_v\end{Smallbmatrix}\GSp(4,\bZ_v)\Bigr)\\
=\omega^{\sfrac{1}{2}}_{\sB}\,\omega^{-\sfrac{1}{2}}_{\sB'}&(\varpi_v)\lambda_{\sB'}\bigl(\GL(2,\bZ_v)\begin{smallbmatrix}\varpi_v\\&1\end{smallbmatrix}\GL(2,\bZ_v)\bigr).
\end{align*}
(It follows from the results in \cite{RobGlobalTheta} that for all classical specializations $\sigma\not\cong\sigma'$ of $\sB,\sB'$ of weights $t>t'\geq 2$ such that they are both discrete series at a finite place $v$ and the product of the central characters is a square, there is a nonzero Yoshida lift of $\sigma\boxtimes\sigma'\otimes (\omega_{\sigma},\omega_{\sigma'})^{\sfrac{1}{2}}\circ\det$ to $\GSp(4)$ with Archimedean component isomorphic to a holomorphic discrete series. This implies the existence of the geometrically irreducible component $\theta(\sB,\sB')\subset \Spec(\bT_{\GSp(4),\ord})$.)

\begin{thm}\label{thm:pstd}
Let $\phi$ be a Dirichlet character unramified outside $S\setminus\{p\}$ with $\phi^2\neq \triv$ and $F_{\theta(\sB,\sB')}={\rm Frac}\left(\bI_{\theta(\sB,\sB')}\right)$. There exists
\[
\resizebox{\displaywidth}{!}{$\displaystyle
\mu^S_{\theta(\sB,\sB'),\phi}\in \Meas\bigl(\bQ^\times\backslash\bA^\times_{\bQ,f}/\hat{\bZ}^{p,\times},\cM_{\GSp(4),\ord}\otimes_{\wt{\Lambda}_{\GSp(4)}}\cM_{\GSp(4),\ord}\bigr)\otimes_{\wt{\Lambda}_{\GSp(4)}} F_{\theta(\sB,\sB')}$}
\]
satisfying the interpolation property: Suppose that $x\in \theta(\sB,\sB')(\bar{\bQ}_p)$ is point where the weight projection map $\wt{\Lambda}_{\GSp(4)}\to \bT_{\GSp(4),\ord}$ is \'{e}tale and has arithmetic image $(l_1,l_2,\xi_1,\xi_2)$ with $l_1\geq l_2\geq 3$. If~$\kappa=(k,\xi)\in\Hom_\cont\bigl(\bQ^\times\backslash\bA^\times_{\bQ,f}/U^p,\bar{\bQ}^\times_p\bigr)$ is an arithmetic point with $l_2\geq k\geq 3$ and $\phi\chi(-1)=(-1)^k$, then\vspace*{-5pt}
\begin{align*}
\mu^S_{\theta(\sB,\sB'),\phi}&(\kappa,x)\\[-3pt]
&\hspace*{-1cm}=2^{-l_1-l_2} \cdot D^S(k-2,\theta(\sB,\sB')_x\times\phi\chi)
\sum_{\varphi\in \sS_{\GSp(4),x}}\hspace*{-3mm}
\frac{\varphi\boxtimes\varphi}{\bfP(\varphi,\varphi)} \\[-3pt]
&\hspace*{-1cm}=2^{-l_1-l_2} \cdot D^S(k-2,\phi\chi)\,D^S\bigl(k-2,\sigma_x\times\sigma'_x\times \omega^{-\sfrac{1}{2}}_{\sigma_x}\omega^{-\sfrac{1}{2}}_{\sigma'_x}\phi\chi\bigr)
\sum_{\varphi\in \sS_{\GSp(4),x}}\hspace*{-3mm}
\frac{\varphi\boxtimes\varphi}{\bfP(\varphi,\varphi)},
\end{align*}
where $\sigma_x,\sigma'_x$ are the (unique) automorphic representations of $\GL(2)$ giving rise to the Hecke eigensystems parameterized by the point in $\sB(\bar{\bQ}_p)\times\sB'(\bar{\bQ}_p)$ induced by $x$ and the natural map $\bI_{\sB}\hat{\otimes}_\cO\bI_{\sB'}\to \bI_{\theta(\sB,\sB')}$, the set $\sS_{\GSp(4),x}$ is an orthogonal basis of the space spanned by ordinary cuspidal holomorphic Siegel modular forms on $\GSp(4)$ of weight $(l_1,l_2)$, tame level $K^p_{\GSp(4)}$ belonging to the Hecke eigenspace parameterized by~$x$, or equivalently of $\theta(\sigma_x,\sigma'_x)$.
\end{thm}

\begin{proof}
In \cite{SLF}, the $p$-adic standard \Lfunctions for ordinary families on symplectic groups (over $\bQ$) are constructed by using the doubling method formula of Garrett \cite{GaPull} and Piatetski-Shapiro--Rallis \cite{PSR}. We~apply that construction to the special case of $\theta(\sB,\sB')$ on $\GSp(4)$. (The Archimedean zeta integral is left as an uncomputed factor in \loccit. It has been calculated in \cite{AZI} and verified to agree with what is expected according to the conjecture of Coates and Perrin--Riou on $p$\nobreakdash-adic \Lfunctions.) The last factor in the formula for $\mu^S_{\theta(\sB,\sB'),\phi}$ here is slightly different from the formula~\cite[(7.0.1)]{SLF} because the interpolation formula is computed by applying $\left<\,,\bar{\varphi}\right>$ to the specialization. If~we apply instead $\bfP(\,,\varphi)$ to the specialization, we~get the above formula.
\end{proof}

\subsection{Comparison of \texorpdfstring{$p$-adic \protect\Lfunctions}{pL}}\label{sec:compare}
Let $\sC_1,\sB,\sB'$ be primitive geometrically irreducible components of $\Spec(\bT_{\GL(2),\ord})$ such that $\Spec(\bT_{\GSp(4),\ord})$ has a geometrically irreducible component $\theta(\sB,\sB')$. By using the results on $p$-adic \Lfunctions in Section~\ref{sec:pL}, we~can deduce the following proposition.

\begin{prop}\label{prop:modpLf}
There exists\vspace*{-5pt}
\begin{multline*}
\cL^{S,\ast}_{\sB,\sC_1},\cL^{S,\ast}_{\sC_1,\sB'}\\
\in\Meas\bigl(\bQ^\times\backslash\bA^\times_{\bQ,f}/U^p,\Lambda_{\GL(2)}\wh{\otimes}_{\cO} \Lambda_{\GSp(4)}\bigr)\otimes_{\wt{\Lambda}_{\GL(2)}\hat{\otimes}_\cO\wt{\Lambda}_{\GSp(4)}} (F_{\sC_1}\hat{\otimes}_\cO F_{\theta(\sB,\sB')})
\end{multline*}
and\vspace*{-5pt}
\[
\cF_{\theta(\sB,\sB')} \in \bigl(\cM_{\GSp(4),\ord}\otimes_{\wt{\Lambda}_{\GSp(4)}} \cM_{\GSp(4),\ord}\bigr)\otimes_{\wt{\Lambda}_{\GSp(4)}} F_{\theta(\sB,\sB')}
\]
satisfying the following interpolation properties: In the setting of Theorem~\ref{thm:pFL1} $\sC_2=\theta(\sB,\sB')$,\vspace*{-5pt}
\begin{align*}
\cL^{S,\ast}_{\sB,\sC_1}(\kappa,x)&=\frac{D^S\bigl(k+\psfrac{l+l_1+l_2}{2},\sigma_x\times\pi_x\times\chi\bigr)}{(-2i)^{l_1+l_2-1}\bfP(h_x,h_x)}\\
\cL^{S,\ast}_{\sC_1,\sB'}(\kappa,x)&=\frac{D^S\bigl(k+\psfrac{l+l_1+l_2}{2},\pi_x\times\sigma'_x\times\omega^{\sfrac{1}{2}}_{\sigma_x}\omega^{-\sfrac{1}{2}}_{\sigma'_x}\chi\bigr)}{(-2i)^{l+1}\bfP(f_x,f_x)}
\end{align*}
and
\begin{align*}
\cF_{\theta(\sB,\sB')}(x)= 2^{-1}i^{-l_1-l_2+1}\,\bfP(h_x,h_x)\sum_{\varphi\in\sS_{\GSp(4),x}} \frac{\varphi\boxtimes\varphi}{\bfP(\varphi,\varphi)}
\end{align*}
with $\pi_x,f_x\in\pi_x$ as in Theorem~\ref{thm:pFL1}, $\sigma_x,\sigma'_x$ as in Theorem~\ref{thm:pstd}, and $h_x\in\sigma_x$ the unique normalized ordinary form fixed by $K^p_{\GL(2)}$. (Note that the weights of the Archimedean components of $\sigma_x,\sigma'_x$ are $l_1+l_2-2,l_1-l_2+2$.)
\end{prop}

\begin{proof}
Applying pullback and change of variable to the Rankin--Selberg $p$-adic \Lfunction in Theorem~\ref{thm:Rankin} shows the existence of $\cL^{S,*}_{\sB,\sC_1},\cL^{S,\ast}_{\sC_1,\sB'}$. In the same way, for a tame Dirichlet character $\phi$ as in Theorem~\ref{thm:pstd}, we~get
\[
\cL^{S,\ast}_{\sB,\sB',\phi}\in \Meas\bigl(\bQ^\times\backslash \bA^\times_{\bQ,f}/\hat{\bZ}^{p,\times},\bI_{\theta(\sB,\sB'}\bigr)\otimes_{\bI_{\theta(\sB,\sB')}} F_{\theta(\sB,\sB')}
\]
satisfying the interpolation property: For all points
\[
(\kappa,x_2)\in \Hom_\cont\bigl(\bQ^\times\backslash \bA^\times_{\bQ,f}/\hat{\bZ}^{p,\times},\bar{\bQ}^\times_p\bigr)\times \theta(\sB,\sB')(\bar{\bQ}_p)
\]
satisfying the conditions in Theorem~\ref{thm:pstd} (where $x_2$ denotes the projection of $x$ to $ \theta(\sB,\sB')(\bar{\bQ}_p)$)
\begin{align*}
\cL^{S,*}_{\sB,\sB',\phi}(\kappa,x)=\frac{D^S\bigl(k-2,\sigma_{x}\times\sigma'_{x}\times \omega^{-\sfrac{1}{2}}_{\sigma_{x}}\omega^{-\sfrac{1}{2}}_{\sigma'_{x}}\chi\phi\bigr)}{(-2i)^{l_1+l_2-1}\bfP(h_{x},h_{x})}.
\end{align*}
A change of variable (\ie $\mu$ to $f\mto \int_{\bQ^\times\backslash\bA^\times_\bQ/\hat{\bZ}^{p,\times}} f(x)x^{-2}\,d\mu(x)$) to the Kubota--Leopoldt $p$-adic \Lfunction recalled in Theorem~\ref{thm:KLp} gives
\[
\cL^{S,\ast}_{\rm KL,\phi}\in \Meas\bigl(\bQ^\times\backslash\bA^\times_\bA/\hat{\bZ}^{p,\times},\cO\bigr)
\]
such that for $\kappa$ as as in Theorem~\ref{thm:pstd}
\[
\cL^{S,\ast}_{\rm KL,\phi}(\kappa)=D^S(k-2,\chi\phi).
\]
The desired $\cF_{\theta(\sB,\sB')}$ can be obtained as $(\cL^{S,\ast}_{\sB,\sB',\phi}\cL^{S,\ast}_{\rm KL,\phi})^{-1} \mu^S_{\theta(\sB,\sB'),\phi}$ with $\mu^S_{\theta(\sB,\sB'),\phi}$ as in Theorem~\ref{thm:pstd}.
\end{proof}

Next, we~show that we can take the $(\bS,\bLam)$-Bessel coefficient on the second factor of $\cF_{\theta(\sB,\sB')}$.

\begin{prop} \label{prop:intBes}
Take nonzero $\cH\in \bI_{\theta(\sB,\sB')}$ such that
\[
\cH \cF_{\theta(\sB,\sB')}\in \cM_{\GSp(4),\ord}\otimes_{\wt{\Lambda}_{\GSp(4)}} \bI_{\theta(\sB,\sB')}.
\]
Given $\bS=\begin{smallbmatrix}\bba&\sfrac{\bbb}{2}\\\sfrac{\bbb}{2}&\bbc\end{smallbmatrix}\in\Sym_2(\bQ)_{>0}$ and $\bLam\in\Hom_\cont(\cK^\times\backslash\bA^\times_{\cK,f},\Lambda^\times_{\GSp(4)})$ extending $\omega_{\theta(\sB,\sB')}=\omega_{\sB}$, where $\cK=\bQ(\sqrt{-\det\bS})$, there exists
\[
\cF_{\theta(\sB,\sB'),\bS,\bLam}
\in \cM_{\GSp(4),\ord}\otimes_{\wt{\Lambda}_{\GSp(4)}} F_{\theta(\sB,\sB')}
\]
such that for $x\in\theta(\sB,\sB')(\bar{\bQ}_p)$ as in Proposition~\ref{prop:modpLf} and not a pole of $\cH$,
\[
\cF_{\theta(\sB,\sB'),\bS,\bLam}(x)=2^{-1}i^{-l_1-l_2+1}\bfP(h_x,h_x)\sum_{\varphi\in\sS_{\GSp(4),x}} \frac{\bes^\dagger_{\bS,\Lambda}(\varphi)\varphi}{\bfP(\varphi,\varphi)}.
\]
\end{prop}
\begin{proof}
We follow the method in \cite[\S10.2]{HsiehYama-Bessel} to construct the desired $\cF_{\theta(\sB,\sB'),\bS,\bLam}$ from the $\cF_{\theta(\sB,\sB')}$ in Proposition~\ref{prop:modpLf}. Fix $c\geq 0$ such that $p^c\alphaS\in \cO_{\cK}$ and an open compact subgroup $U^p_\cK\subset\bA^{\times,p}_\cK$ such that $\bLam$ factors through the quotient by $U^p_\cK$. Given a positive integer $n$, we~let
\begin{align*}
U_{\cK,p,n}=\bZ^\times_p (1+p^{n+c}\bZ_p\alphaS),
\qquad U_{\cK,n}=U^p_\cK U_{\cK,p,n},
\end{align*}
and
\[
\rho_n:\wt{\Lambda}_{\GSp(4)}\to \cO\big[\,T^1_{\GSp(4)}(\bZ/p^n\bZ_)\,\big]
\]
be the natural projection induced by $T^1_{\GSp(4)}(\bZ_p)\to T^1_{\GSp(4)}(\bZ/p^n\bZ)$. Put
\begin{align*}
\bI_{\theta(\sB,\sB'),n}&=\bI_{\theta(\sB,\sB')}\otimes_{\wt{\Lambda}_{\GSp(4)},\rho_n} \cO\big[\,T^1_{\GSp(4)}(\bZ/p^n\bZ)\,\big].
\end{align*}
Then $\rho_n$ naturally induces $\rho_n:\bI_{\theta(\sB,\sB')}\to \bI_{\theta(\sB,\sB'),n}$. Taking the $q$-expansion of the second factor at $\begin{smallbmatrix}A\\&D\end{smallbmatrix}\in \GSp(4,\bA_{\bQ,f})$ and taking the coefficient indexed by $\bS$ gives a map
\[
\varepsilon_{\qexp,\bS}\left(\,\cdot\,, \begin{smallbmatrix}A\\&D\end{smallbmatrix}\right):\cM_{\GSp(4),\ord}\otimes_{\wt{\Lambda}_{\GSp(4)}}\cM_{\GSp(4),\ord}\to \cM_{\GSp(4)}.
\]
Define $\Theta_n\in \cM_{\GSp(4),\ord}\otimes_{\wt{\Lambda}_{\GSp(4)}} \bI_{\theta(\sB,\sB'),n}$ as
\[
\resizebox{\displaywidth}{!}{$\displaystyle
\sum_{\fz\in \cK^\times\bA^\times_{\bQ,f}\backslash\bA^\times_{\cK,f}/U_{\cK,n}} \hspace{-1em}\rho_n\biggl(\lambda^{-n-c}\cdot\bLam(\fz)^{-1} \,\varepsilon_{\qexp,\bS}\biggl(\cH \cF_{\theta(\sB,\sB')},\begin{smallbmatrix}\imath_\bS(\fz)\\&\ltrans{\imath_\bS(\bar{\fz})}\end{smallbmatrix}\begin{Smallbmatrix}p^{2}\\&p\\&&1\\&&&p\end{Smallbmatrix}^{n+c}_p\biggr)\biggr)[\fz],$}
\]
with $\lambda$ equal to the product of $\omega_{\sB,p}(p)$ and the eigenvalue of the $\bU_p$-operator associated to $\begin{Smallbmatrix}p\\&1\\&&p^{-1}\\&&&1\end{Smallbmatrix}$ corresponding to $\theta(\sB,\sB')$.
For $y\in\bZ_p$, we~have
\begin{multline*}
\imath_\bS(1+p^{n+c}y\alphaS)\\
=\begin{bsm}1&p^{n+c}y\\&1\\ &&1\\&&-p^{n+c}y&1\end{bsm}\begin{bsm}(1+p^{n+c}y\alphaS)(1+p^{n+c}y\alphabS)\\-p^{n+c}y\alphaS\alphabS&1\\ &&1&p^{n+c}y\alphaS\alphabS\\ && &(1+p^{n+c}y\alphaS)(1+p^{n+c}y\alphabS)\end{bsm}
\end{multline*}
and
\[
\resizebox{\displaywidth}{!}{$\displaystyle
\begin{aligned}
\rho_n\biggl(&\varepsilon_{\qexp,\bS}\biggl(\sum_{y\in\bZ/p} \cH\cF_{\theta(\sB,\sB')},\begin{Smallbmatrix}\imath_\bS(\fz(1+p^{n+c}\alphaS y))\\&\ltrans{\imath_\bS}(\bar{\fz}(1+p^{n+c}\alphabS y))\end{Smallbmatrix}\begin{Smallbmatrix}p^{2}\\&p\\&&1\\&&&p\end{Smallbmatrix}^{n+1+c}_p\biggr)\biggr)\\
&=\rho_n\biggl(\varepsilon_{\qexp,\bS}\biggl(\sum_{y\in\bZ/p} \cH\cF_{\theta(\sB,\sB')},\begin{Smallbmatrix}\imath_\bS(\fz)\\&\ltrans{\imath_\bS(\bar{\fz})}\end{Smallbmatrix}\begin{Smallbmatrix}1&p^{n+c}y\\&1\\ &&1\\&&-p^{n+c}y&1\end{Smallbmatrix}\begin{Smallbmatrix}p^{2}\\&p\\&&1\\&&&p\end{Smallbmatrix}^{n+1+c}_p\biggr)\biggr)\\
&=\rho_n\biggl(\varepsilon_{\qexp,\bS}\biggl(\sum_{y\in\bZ/p} \cH\cF_{\theta(\sB,\sB')},\begin{Smallbmatrix}\imath_\bS(\fz)\\&\ltrans{\imath_\bS(\bar{\fz})}\end{Smallbmatrix}\begin{Smallbmatrix}p^{2}\\&p\\&&1\\&&&p\end{Smallbmatrix}^{n+c}_p\begin{Smallbmatrix}1&y\\&1\\ &&1\\&&-y&1\end{Smallbmatrix}\begin{Smallbmatrix}p^{2}\\&p\\&&1\\&&&p\end{Smallbmatrix}_p\biggr)\biggr)\\
&=\rho_n\biggl(\lambda_{\GSp(4),2,1}\cdot \varepsilon_{\qexp,\bS}\biggl( \cH\cF_{\theta(\sB,\sB')},\begin{Smallbmatrix}\imath_\bS(\fz)\\&\ltrans{\imath_\bS(\bar{\fz})}\end{Smallbmatrix}\begin{Smallbmatrix}p^{2}\\&p\\&&1\\&&&p\end{Smallbmatrix}^{n+c}_p\biggr)\biggr),
\end{aligned}$}\]
from which it follows that
\[
\rho_n(\Theta_{n+1})=\Theta_n.
\]
Therefore, the $\Theta_n$'s define an element $\Theta\in \cM_{\GSp(4),\ord}\otimes_{\wt{\Lambda}_{\GSp(4)}}\bI_{\theta(\sB,\sB')}$. Notice that for ordinary $\varphi$ invariant under $\{g\in\GSp(4,\bZ_p):g\bmod p^n \in U_{\GSp(4)}(\bZ/p^n)\}$, $\bes^\dagger_{\bS,\Lambda}(\varphi)$ can be computed with $m_1=n+c,m_2=0$ in \eqref{eq:bes-dagger}. Then from the definition of $\Theta_n$'s, we~see that up to a scalar, $\cH^{-1}\Theta$ gives the desired $\cF_{\theta(\sB,\sB'),\bS,\bLam}$.
\end{proof}

\begin{prop}\label{prop:compare}
Suppose that $\sC_1$ is primitive and $\bm{f}\in \cM_{\GL(2),\ord}$ is the Hida family corresponding to $\sC_1$ normalized such that the first Fourier coefficient is $1$. Then
\begin{equation}\label{eq:compare}
\mu^S_{\sC_1,\theta(\sB,\sB')}
=\bbc\sqrt{\det\bS}\,2^{3}i^{-1} {\bm f}_{\bbc}{\bm f}\cdot \cL^{S,\ast}_{\sB,\sC_1}\cL^{S,\ast}_{\sC_1,\sB'} \cdot \cF_{\theta(\sB,\sB'),\bS,\bLam}
\end{equation}
\end{prop}
\begin{proof}
It suffices to check that in the setting of Theorem~\ref{thm:pFL1} with $\sC_2=\theta(\sB,\sB')$, the evaluations of both sides agree at all $(\kappa,x)$ with $l_1=l_2=l$. It follows from Propositions~\ref{prop:modpLf}, \ref{prop:intBes} that
\begin{multline*}
\text{RHS}(\kappa,x)=\bbc\sqrt{\det\bS}\,2^{3}i^{-1} f_{x,\bbc}f_x\cdot 2^{-3l-1}i^{l+1}\\
\times\frac{D^S\left(k+\sfrac{3l}{2},\sigma_x\times\pi_x\times\chi\right) D^S\Bigl(k+\sfrac{3l}{2},\pi_x\times\sigma'_x\times\omega^{\sfrac{1}{2}}_{\sigma_x}\omega^{-\sfrac{1}{2}}_{\sigma'_x}\times\chi\Bigr)}{\bfP(f_x,f_x)}
\\[-8pt]
\times\sum_{\varphi\in \sS_{\GSp(4),x}}\hspace{-3mm}\frac{\bes^\dagger_{\bS,\Lambda}(\varphi)\varphi}{\bfP(\varphi,\varphi)},
\end{multline*}
where $f_x$ denotes the specialization of ${\bm f}$ at $x$. Note that when $\Pi_x=\theta(\sB,\sB')_x$, the \Lfunction for $\Pi_x$ decomposes as the product of \Lfunctions for $\sigma_x$ and \hbox{$\sigma'_x\otimes \omega^{\sfrac{1}{2}}_{\sigma_x}\omega^{-\sfrac{1}{2}}_{\sigma'_x}\circ\det$}, and we have
\[
D^S(s,\Pi_x\times\pi_x\times\chi)=D^S\left(s,\sigma_x\times\pi_x\times\chi\right) D^S\bigl(s,\pi_x\times\sigma'_x\times\omega^{\sfrac{1}{2}}_{\sigma_x}\omega^{-\sfrac{1}{2}}_{\sigma'_x}\times\chi\bigr).
\]
Hence,
\begin{multline*}
\text{RHS}(\kappa,x)\\
=\bbc\sqrt{\det\bS}\, 2^{-3l+2} i^{l} \cdot D^S\Bigl(k+\frac{3l}{2},\Pi_x\times\pi_x\times\chi\Bigr)\cdot \frac{f_{x,\bbc}f_x}{\bfP(f_x,f_x)}\sum_{\varphi\in \sS_{\GSp(4),x}}\frac{\bes^\dagger_{\bS,\Lambda}(\varphi)\varphi}{\bfP(\varphi,\varphi)}
\end{multline*}
which equals exactly $\mu^S_{\sC_1,\theta(\sB,\sB')}$ by the formula in Theorem~\ref{thm:pFL1}.
\end{proof}

\Subsection{The four-variable \texorpdfstring{$p$-adic \protect\Lfunction}{pL} and its interpolation formula II}
With Proposition~\ref{prop:compare}, we~can deduce a formula for the Archimedean zeta $I_\infty(k,\cD_{l_1,l_2},\cD_l,\Lambda_\infty)$ appearing in the interpolation formula in Theorem~\ref{thm:pFL1} and finish the proof of Theorem~\ref{thm:main}.

\begin{proof}[Proof of Theorem~\ref{thm:main}.]
We choose $K^p_{\GL(2)},K^p_{\GSp(4)}$ and primitive Hida families $\sB,\sB'$ of tame level $K^p_{\GL(2)}$ such that $\Spec(\bT_{\GSp(4),\ord})$ has an irreducible component $\theta(\sB,\sB')$. With such a choice, we~can further choose $\bS$ and $\bLam$ such that $\cF_{\theta(\sB,\sB'),\bS,\bLam}\neq 0$. (By the interpolation property of $\cF_{\theta(\sB,\sB'),\bS,\bLam}$ in Proposition~\ref{prop:intBes}, to~show the existence of such a $\bS$ and $\bLam$, it suffices to show that there exists $x$ satisfying the conditions there for which $B^\dagger_{\bS,\Lambda}(\varphi)\neq 0$ for some $\varphi\in \sS_{\GSp(4),x}$. Take an $x$ with corresponding weight $(l_1,l_2)$, $l_1\gg l_2\gg 0$ and $\varphi\in \sS_{\GSp(4),x}$. One can choose~$\bS,\Lambda$ such that the usual Bessel period $B_{\bS,\Lambda}(\varphi)\neq 0$. Then by \cite[Prop.\,2.7.1]{pFL}, we~know that $B^\dagger_{\bS,\Lambda}(\varphi)\neq 0$.) Then we can choose the finite set $S$ such that the conditions in Section~\ref{sec:setup} is satisfied. By the primitivity of $\sC_1$ and our assumption on $K^p_{\GL(2)}$ made at the beginning of Section~\ref{sec:pL}, $\bm{f}_\bbc\neq 0$. Hence, both sides of the identity in Proposition~\ref{prop:compare} are nonzero elements in $\cM_{\GSp(4),\ord}\otimes_{\wt{\Lambda}_{\GSp(4)}} F_{\theta(\sB,\sB')}$. (There are many interpolation points corresponding to $s$ belonging to the absolute convergence range, at which one can check that the evaluations are nonzero.)

Let $(\kappa,x)$ be a point of $\Hom_\cont\bigl(\bQ^\times\backslash\bA^\times_{\bQ,f}/U^p,\bar{\bQ}^\times_p \bigl)\times\sC_1(\bar{\bQ}_p)\times \theta(\sB,\sB')(\bar{\bQ}_p)$ as in Theorem~\ref{thm:pFL1} with $\sC_2=\theta(\sB,\sB')$. Then, for the evaluations of both sides of \eqref{eq:compare} at $(\kappa,x)$, we~have\vspace*{-5pt}
\begin{multline*}
\text{RHS}(\kappa,x)
=\bbc\sqrt{\det\bS}\,2^{3}i^{-1} f_{x,\bbc}f_x\cdot 2^{-l-l_1-l_2-1}i^{l+1}
\sum_{\varphi\in \sS_{\GSp(4),x}}\hspace*{-5mm}\frac{\bes^\dagger_{\bS,\Lambda}(\varphi)\varphi}{\bfP(\varphi,\varphi)} \\
\times\frac{D^S\left(k+\frac{l+l_2+l_2}{2},\sigma_x\times\pi_x\times\chi\right) D^S\Bigl(k+\frac{l+l_1+l_2}{2},\pi_x\times\sigma'_x\times\omega^{\sfrac{1}{2}}_{\sigma_x}\omega^{-\sfrac{1}{2}}_{\sigma'_x}\times\chi\Bigr)}{\bfP(f_x,f_x)}
\end{multline*}
by the interpolation properties of $\cL^{S,\ast}_{\sB,\sC_1},\cL^{S,\ast}_{\sC_1,\sB'}, \cF_{\theta(\sB,\sB'),\bS,\bLam}$, and\vspace*{-5pt}
\begin{multline*}
{\rm LHS}(\kappa,x)
= i^{r_{\sLambda,1}-r_{\sLambda,2}} f_{x,\bbc}f_x \cdot\frac{I_\infty(k,\cD_{l_1,l_2},\cD_l,\Lambda_\infty)}{E_\infty\Bigl(k+\frac{l+l_1+l_2}{2},\cD_{l_1,l_2}\times\cD_l\times\chi\Bigr)}\sum_{\varphi\in \sS_{\GSp(4),x}}\hspace*{-5mm}\frac{\bes^\dagger_{\bS,\Lambda}(\varphi)\varphi}{\bfP(\varphi,\varphi)} \\
\times\frac{D^S\left(k+\frac{l+l_2+l_2}{2},\sigma_x\times\pi_x\times\chi\right) D^S\Bigl(k+\frac{l+l_1+l_2}{2},\pi_x\times\sigma'_x\times\omega^{\sfrac{1}{2}}_{\sigma_x}\omega^{-\sfrac{1}{2}}_{\sigma'_x}\times\chi\Bigr)}{\bfP(f_x,f_x)}
\end{multline*}
by Theorem~\ref{thm:pFL1}. It follows that
\begin{multline}\label{eq:arch-zeta}
I_\infty(k,\cD_{l_1,l_2},\cD_l,\Lambda_\infty)\\
= \bbc\sqrt{\det\bS} \,2^{-l-l_1-l_2+2}i^{l-r_{\sLambda,1}+r_{\sLambda,2}} \cdot E_\infty\Bigl(k+\psfrac{l+l_1+l_2}{2},\cD_{l_1,l_2}\times\cD_l\times\chi\Bigr),
\end{multline}
for all $(k,l,l_1,l_2)$ which equals the algebraic part of the projection of an arithmetic $(\kappa,x)$ to the weight space such that $x$ is not a pole of either side of\eqref{eq:compare} and ${\rm RHS}(\kappa,x)$ and ${\rm LHS}(\kappa,x)$ are nonzero.

Since we have made the choices such that both sides of \eqref{eq:compare} are nonzero, the points for which the weight projection map \eqref{eq:wtproj} is not \'{e}tale at $x$ or ${\rm LHS}(\kappa,x)={\rm LHS}(\kappa,x)=0$ are not Zariski dense. For any $(k,l,l_1,l_2)$ satisfying \eqref{eq:kl}, the classical points $(\kappa,x)$ whose projections to the weight space has algebraic part equal to $(k,l,l_1,l_2)$ are Zariski dense, so there exist $(\kappa,x)$ which satisfies the conditions for the above comparison to deduce \eqref{eq:arch-zeta} for the given $(k,l,l_1,l_2)$. Thus, \eqref{eq:arch-zeta} is true for all $(k,l,l_1,l_2)$ satisfying \eqref{eq:kl}. Plugging it into the interpolation formula in Theorem~\ref{thm:pFL1} shows that
\[
\cL^S_{\sC_1,\sC_2,\beta_1,\beta_2}=\bigl(\bbc\sqrt{\det\bS}\bigr)^{-1}\,2^{-2}\cdot \varepsilon_{\qexp,\beta_1,\beta_2}\left(\mu^S_{\sC_1,\sC_2}\right)
\]
is the desired $p$-adic \Lfunction.
\end{proof}
\begin{thebibliography}{BGR23}
\bibitem[Liu23]{pFL} Z. Liu, “p-adic L-functions for GSp(4) × GL(2)”, 2023, arXiv:2308.08533.
\bibitem[Hid02]{HidaCon} H. Hida, “Control theorems of coherent sheaves on Shimura varieties of PEL type”, J. Inst. Math. Jussieu 1 (2002), no. 1, p. 1–76.
\bibitem[Liu20]{SLF} Z. Liu, “p-adic L-functions for ordinary families on symplectic groups”, J. Inst. Math. Jussieu 19 (2020), no. 4, p. 1287–1347.
\bibitem[Fur93]{Furusawa} M. Furusawa – “On L-functions for GSp(4) × GL(2) and their special values”, J. reine angew. Math. 438 (1993), p. 187–218.
\bibitem[Gar89]{GaKl} P. B. Garrett, “Integral representations of Eisenstein series and L-functions”, in Number theory, trace formulas and discrete groups (Oslo,1987), Academic Press, Boston, MA, 1989, p. 241–264.
\bibitem[Mor14]{Morimoto-GSp4GL2-I} K. Morimoto – “On L-functions for quaternion unitary groups of degree 2 and GL(2) (with an appendix by M. Furusawa and A. Ichino)”, Internat. Math. Res. Notices (2014), no. 7, p. 1729–1832.
\bibitem[Mor18]{Morimoto-GSp4GL2-II} K. Morimoto, “On tensor product L-functions for quaternion unitary groups and GL(2) over totally real number fields: mixed weight cases”, Adv. Math. 337 (2018), p. 317–362.
\bibitem[Shi00]{Sh00} G. Shimura – Arithmeticity in the theory of automorphic forms, Math. Surveys and Monographs, vol. 82, American Mathematical Society, Providence, RI, 2000.
\bibitem[Liu19]{NHF} Z. Liu – “Nearly overconvergent Siegel modular forms”, Ann. Inst. Fourier (Grenoble) 69 (2019), no. 6, p. 2439–2506.
\bibitem[FC90]{FC} G. Faltings \& C.-L. Chai – Degeneration of abelian varieties, Ergeb. Math. Grenzgeb. (3), vol. 22, Springer-Verlag, Berlin, 1990.
\bibitem[Eis12]{EisDiff} E. E. Eischen – “p-adic differential operators on automorphic forms on unitary groups”, Ann. Inst. Fourier (Grenoble) 62 (2012), no. 1, p. 177–243.
\bibitem[Hid04]{HidaPAF} H. Hida, p-adic automorphic forms on Shimura varieties, Springer Monographs in Math., Springer-Verlag, New York, 2004.
\bibitem[Kat73]{TDwork} N. Katz – “Travaux de Dwork”, in Séminaire Bourbaki (1971/1972), Lect. Notes in Math., vol. 317, Springer, Berlin, 1973, Exp. No. 409, p. 167–200.
\bibitem[CPR89]{CoPerrin} J. Coates \& B. Perrin-Riou – “On p-adic L-functions attached to motives over Q”, in Algebraic number theory, Adv. Stud. Pure Math., vol. 17, Academic Press, Boston, MA, 1989, p. 23–54.
\bibitem[Coa91]{Coates} J. Coates – “Motivic p-adic L-functions”, in L-functions and arithmetic (Durham, 1989), London Math. Soc. Lecture Note Ser., vol. 153, Cambridge University Press, Cambridge, 1991, p. 141–172.
\bibitem[Hid88]{Hi88} H. Hida – “A p-adic measure attached to the zeta functions associated with two elliptic modular forms. II”, Ann. Inst. Fourier (Grenoble) 38 (1988), no. 3, p. 1–83.
\bibitem[CH20]{Chen-Hsieh} S.-Y. Chen \& M.-L. Hsieh – “On primitive p-adic Rankin-Selberg L-functions”, in Development of Iwasawa theory–the centennial of K. Iwasawa’s birth, Adv. Stud. Pure Math., vol. 86, Mathematical Society of Japan, Tokyo, 2020, p. 195–242.
\bibitem[Hsi21]{HsiehTriple} M.-L. Hsieh – “Hida families and p-adic triple product L-functions”, Amer. J. Math. 143 (2021), no. 2, p. 411–532.
\bibitem[Rob01]{RobGlobalTheta} B. Roberts – “Global L-packets for GSp(2) and theta lifts”, Doc. Math. 6 (2001), p. 247– 314.
\bibitem[Gar84]{GaPull} P. B. Garrett – “Pullbacks of Eisenstein series; applications”, in Automorphic forms of several variables (Katata, 1983), Progress in Math., vol. 46, Birkhäuser Boston, Boston, MA, 1984, p. 114–137.
\bibitem[PSR87]{PSR} I. Piatetski-Shapiro \& S. Rallis – “L-functions for the classical groups”, in Explicit constructions of automorphic L-functions, Lect. Notes in Math., vol. 1254, Springer-Verlag, Berlin, 1987, p. 1–52.
\bibitem[Liu21]{AZI} Z. Liu, “The doubling archimedean zeta integrals for p-adic interpolation”, Math. Res. Lett. 28 (2021), no. 1, p. 145–173.
\bibitem[HY24]{HsiehYama-Bessel} M.-L. Hsieh \& S. Yamana – “Bessel periods and anticyclotomic p-adic spinor L-functions”, Trans. Amer. Math. Soc. 377 (2024), no. 8, p. 5617–5672.
\end{thebibliography}
\end{document}
